% P11 paper module: canonical O3j reconciliation for the Dirichlet--Riesz reduction.

\subsection{Dirichlet--Riesz reduction and the Gamma-action reconciliation}\label{subsec:o3j-reconciled}

Keep $0<R<S<T_0$ fixed and let
$h\in C_c^\infty((-S,S))$ be the smooth odd annulus vector used above, with support
disjoint from $[-R,R]$.  Define
\[
 u_h:=G_{R,T_0}^{-1}J_{R,S}^*G_{S,T_0}h,
 \qquad
 g_h=h-J_{R,S}u_h.
\tag{OJ.1}\label{eq:o3j-uh}
\]
Then $J_{R,T_0}u_h$ is the $T_0$-metric orthogonal projection of
$J_{S,T_0}h$ onto $\Ran J_{R,T_0}$.  On the underlying $L^2$ spaces this means
\[
 q_{T_0}^X(E_{R,T_0}u_h,E_{R,T_0}v)
 =q_{T_0}^X(E_{S,T_0}h,E_{R,T_0}v)
 \qquad(v\in\KXR{R}).
\tag{OJ.2}\label{eq:o3j-variational}
\]
Here $E_{R,T_0}$ and $E_{S,T_0}$ are the ordinary $L^2$ zero-extension maps; this
notation is kept distinct from the graph maps $J_{R,S}$.

Let
\[
 \Sigma_R^{[T_0]}
 :=E_{R,T_0}^*\Sigma_{T_0}E_{R,T_0}.
\tag{OJ.3}\label{eq:o3j-schur-restriction}
\]
It is bounded, positive and selfadjoint on $L^2(-R,R)$.  Equation
\eqref{eq:o3j-variational} is therefore the weak equation
\[
 (C_{\Gamma,R}+\Sigma_R^{[T_0]})u_h=r_h,
 \qquad
 r_h=r_{\Gamma,h}+r_{\sigma,h},
\tag{OJ.4}\label{eq:o3j-forcing-equation}
\]
where
\[
 r_{\Gamma,h}
 :=E_{R,T_0}^*C_{\Gamma,T_0}E_{S,T_0}h,
 \qquad
 r_{\sigma,h}
 :=E_{R,T_0}^*\Sigma_{T_0}E_{S,T_0}h.
\tag{OJ.5}\label{eq:o3j-forcing-split}
\]
The second term belongs to $L^2(-R,R)$ by boundedness.  The first term requires a
careful operator-domain argument; it is here that the earlier shortcut is replaced.

\begin{lemma}[Gamma action on the smooth interior core]\label{lem:gamma-smooth-action}
Let $T>0$ and $\phi\in C_c^\infty((-T,T))$.  Put
\[
 \boxed{
 \mathcal G_\phi
 :=\mathcal F^{-1}\!\left(m_\Gamma\widehat{E_T\phi}\right).
 }
\tag{OJ.6}\label{eq:gamma-Gphi}
\]
Then
\[
 \mathcal G_\phi\in H^N(\mathbb R)
 \qquad\text{for every finite }N,
\tag{OJ.7}\label{eq:gamma-Gphi-sobolev}
\]
and
\[
 \boxed{
 \phi\in\mathcal D(C_{\Gamma,T}),
 \qquad
 C_{\Gamma,T}\phi=P_T\mathcal G_\phi.
 }
\tag{OJ.8}\label{eq:gamma-smooth-action}
\]
\end{lemma}

\begin{proof}
Since $\phi$ is smooth and compactly supported strictly inside $(-T,T)$,
$\widehat{E_T\phi}$ is Schwartz.  The multiplier
$m_\Gamma(\xi)\asymp\log(2+|\xi|)$ has only logarithmic growth, so
$m_\Gamma\widehat{E_T\phi}$ has every finite polynomial $L^2$ moment; this proves
\eqref{eq:gamma-Gphi-sobolev}.  For every Gamma-form vector $v$,
Plancherel gives
\[
 \mathfrak c_{\Gamma,T}[\phi,v]
 =\langle P_T\mathcal G_\phi,v\rangle_{L^2(-T,T)}.
\]
The representation theorem for the closed positive form therefore puts this particular
$\phi$ in the operator domain and identifies the operator action as in
\eqref{eq:gamma-smooth-action}.  No general assertion that every smooth vector belongs
to the operator domain of an arbitrary form-defined operator is being used.
\end{proof}

Apply Lemma~\ref{lem:gamma-smooth-action} to
\[
 \phi:=E_{S,T_0}h\in C_c^\infty((-T_0,T_0)).
\]
Then
\[
 r_{\Gamma,h}=P_R\mathcal G_\phi.
\tag{OJ.9}\label{eq:o3j-gamma-forcing}
\]
The restriction is arbitrarily Sobolev-regular on the fixed interval $(-R,R)$.  After
zero extension at $\pm R$ there may be fixed boundary jumps, but such jumps have Fourier
decay $O(|\xi|^{-1})$; consequently
\[
 E_Rr_{\Gamma,h}
 \in\bigcap_{\alpha<\infty}\HG^\alpha(\mathbb R).
\tag{OJ.10}\label{eq:o3j-gamma-logregular}
\]
In particular $r_{\Gamma,h}\in L^2(-R,R)$, and hence $r_h\in L^2(-R,R)$.

Since $\Sigma_R^{[T_0]}$ is a bounded selfadjoint perturbation of $C_{\Gamma,R}$,
\[
 \mathcal D(C_{\Gamma,R}+\Sigma_R^{[T_0]})
 =\mathcal D(C_{\Gamma,R}).
\]
Solving~\eqref{eq:o3j-forcing-equation} therefore gives the genuine operator-domain
gain
\[
 \boxed{
 u_h
 =(C_{\Gamma,R}+\Sigma_R^{[T_0]})^{-1}r_h
 \in\mathcal D(C_{\Gamma,R}).
 }
\tag{OJ.11}\label{eq:o3j-domain-gain}
\]

\begin{remark}[O3j local firewall and status]\label{rem:o3j-firewall}
The argument through~\eqref{eq:o3j-domain-gain} by itself proves neither
\[
 E_Ru_h\in\HG^{m_h+3/2}
\]
nor higher logarithmic regularity of the Schur forcing $r_{\sigma,h}$.  In particular,
one must not replace $G_{R,T_0}^{-1}$ by the global Fourier multiplier
$m_\Gamma^{-1}$; the finite-window metric inverse remains a genuine Dirichlet/Riesz
operator.  The next subsection closes this local regularity gap by a separate
fixed-window Sobolev/interpolation argument, without making that replacement.  Thus
Open Problem~\ref{open:log}, stated immediately before this reconciliation block, is a
historical formulation of the gate and is resolved positively by
Theorem~\ref{thm:o3k-complement} below.
\end{remark}

% P11 paper module: R12 positive resolution of the logarithmic complement gate.

\subsection{Positive Sobolev bootstrap for the explicit complement}\label{subsec:o3k-log-complement}

The previous subsection deliberately stopped at the operator-domain gain
\eqref{eq:o3j-domain-gain}.  At fixed terminal horizon, however, the source cutoff has
additional structure: after resolving the finite martingale coordinates it is a finite
family of interval cutoffs and translations.  This permits a small positive Sobolev
bootstrap.  Since any positive Sobolev order dominates every finite logarithmic order,
this closes the logarithmic complement gate.

For $s\in\mathbb R$ let
\[
 X_T^s:=\{F\in H^s(\mathbb R):\operatorname{supp}F\subset[-T,T]\}
\]
with the inherited $H^s$ norm.  We use the following standard cutoff fact.

\begin{lemma}[Small Sobolev cutoff range]\label{lem:o3k-cutoff}
For every bounded interval $I\subset\mathbb R$, multiplication by $1_I$ is bounded on
$H^s(\mathbb R)$ whenever $|s|<1/2$.  More generally, for fixed finite $a\ge0$ it is
bounded on the weighted Fourier space
\[
 \mathcal H^{s,a}
 :=\left\{F:\int_{\mathbb R}
 \langle\xi\rangle^{2s}[\log(2+|\xi|)]^{2a}
 |\widehat F(\xi)|^2\,d\xi<\infty\right\}
\]
whenever $|s|<1/2$.
\end{lemma}

\begin{proof}
The first statement is the classical fractional-Sobolev interval-multiplier theorem.
For the logarithmically refined statement put
\[
 w_{s,a}(\xi):=\langle\xi\rangle^{2s}[\log(2+|\xi|)]^{2a}.
\]
For $|s|<1/2$, $w_{s,a}$ is an $A_2$ weight: a direct interval-average estimate reduces
at large scale to the power range $-1<2s<1$, while the logarithmic factor is slowly
varying and does not alter the open $A_2$ range.  Multiplication by a half-line in
physical space becomes, after Fourier transform, a fixed linear combination of the
identity and the Hilbert transform.  The weighted Hilbert-transform theorem on
$L^2(w_{s,a})$ therefore gives boundedness for half-line cutoffs; differences of two
translated half-lines give bounded intervals.
\end{proof}

\begin{lemma}[Fixed-window Schur Sobolev window]\label{lem:o3k-schur-window}
For every fixed $T>0$ there exists $\varepsilon_T>0$ such that
\[
 B_T=(I+R_T^*R_T)^{-1}:X_T^s\longrightarrow X_T^s
\]
is bounded for all $|s|<\varepsilon_T$.  Consequently $H_T$, $H_T^*$ and
$\Sigma_T=H_TB_TH_T^*$ are bounded on $X_T^s$ after possibly shrinking
$\varepsilon_T$.
\end{lemma}

\begin{proof}
At fixed $T$ the sum defining $R_T$ is effectively finite.  By
\eqref{eq:projected-mark}, each martingale coordinate contains a scalar factor
\[
 1_{\{j<J_{p,T}(u)\}},
\]
which, up to endpoints, is the indicator of
\[
 \left\{|u|<T-\frac{j+1}{2}\log p\right\}.
\]
Thus, after resolving the finite martingale coordinates, $R_T$ is a finite sum of
compositions of fixed translations, interval cutoffs, zero extension/restriction and
fixed finite-dimensional coefficient maps.  Lemma~\ref{lem:o3k-cutoff} therefore gives
bounded maps
\[
 R_T:X_T^s\longrightarrow Z_T^s
 \qquad(|s|<1/2)
\]
on the corresponding finite vector-valued Sobolev target.  Applying the same statement
at exponent $-s$ and dualizing gives
\[
 R_T^*:Z_T^s\longrightarrow X_T^s.
\]
Hence
\[
 A_T^{\rm rest}:=I+R_T^*R_T
\]
is bounded on $X_T^s$ for $|s|<1/2$.

At $s=0$, $A_T^{\rm rest}\ge I$, so it is an isomorphism on $X_T^0=L^2(-T,T)$.
The spaces $X_T^s$ form a complex interpolation scale for $|s|<1/2$: they are the
ranges of the common interval-cutoff projection, which is bounded at the endpoint
exponents by Lemma~\ref{lem:o3k-cutoff}.  The local stability theorem for invertible
operators on complex interpolation scales (the \v Sne\u\i berg theorem) now yields an
$\varepsilon_T>0$ such that $A_T^{\rm rest}$ remains invertible on $X_T^s$ whenever
$|s|<\varepsilon_T$.  This is the asserted boundedness of $B_T$.

The hub $H_T$ is itself a finite sum of fixed translation differences and interval
extension/restriction, hence is bounded on the same small Sobolev scale.  Boundedness of
$H_T^*$ follows by using the negative exponents and duality, and the assertion for
$\Sigma_T$ follows by composition.
\end{proof}

\begin{proposition}[Positive Sobolev regularity of the O3j forcing]\label{prop:o3k-forcing}
For the fixed triple $0<R<S<T_0$ and smooth annulus vector $h$ of
Subsection~\ref{subsec:o3j-reconciled}, there exists $s_1>0$ such that
\[
 \boxed{E_Rr_h\in H^{s_1}(\mathbb R).}
\tag{OJ.12}\label{eq:o3k-forcing-sobolev}
\]
\end{proposition}

\begin{proof}
Put $\phi=E_{S,T_0}h\in C_c^\infty((-T_0,T_0))$.  Choose
$0<s_1<\min\{\varepsilon_{T_0},1/2\}$ in
Lemma~\ref{lem:o3k-schur-window}.  Since $\phi$ is smooth and supported strictly inside
$(-T_0,T_0)$,
\[
 H_{T_0}^*\phi\in X_{T_0}^{s_1}.
\]
Therefore
\[
 \Sigma_{T_0}\phi
 =H_{T_0}B_{T_0}H_{T_0}^*\phi
 \in X_{T_0}^{s_1}.
\]
Restriction to $(-R,R)$ and zero extension back to the line preserve $H^{s_1}$, after
shrinking $s_1$ if necessary, so
\[
 E_Rr_{\sigma,h}\in H^{s_1}(\mathbb R).
\]
On the Gamma side, Lemma~\ref{lem:gamma-smooth-action} gives
$r_{\Gamma,h}=P_R\mathcal G_\phi$ with
$\mathcal G_\phi\in H^N(\mathbb R)$ for every finite $N$.  Lemma~\ref{lem:o3k-cutoff}
therefore gives $E_Rr_{\Gamma,h}\in H^{s_1}$ for every sufficiently small
$s_1<1/2$.  Adding the two terms proves~\eqref{eq:o3k-forcing-sobolev}.
\end{proof}

We next bootstrap through the genuine finite-window form operator, rather than replacing
it by a full-line Fourier multiplier.  Put $m=m_\Gamma$ and for $|s|<1/2$ define
\[
 \mathcal V^s
 :=\left\{F:\int_{\mathbb R}
 \langle\xi\rangle^{2s}m(\xi)|\widehat F(\xi)|^2\,d\xi<\infty\right\},
\qquad
 V_R^s:=\{F\in\mathcal V^s:\operatorname{supp}F\subset[-R,R]\}.
\tag{OJ.13}\label{eq:o3k-Vs}
\]
Since $m(\xi)\asymp\log(2+|\xi|)$, Lemma~\ref{lem:o3k-cutoff} applies to these weights.
The $V_R^s$ are therefore complemented support subspaces of the full weighted Fourier
scale and satisfy the complex interpolation law
\[
 [V_R^{-\delta},V_R^{\delta}]_\theta
 =V_R^{(2\theta-1)\delta}
\tag{OJ.14}\label{eq:o3k-interpolation}
\]
for every sufficiently small fixed $\delta>0$.  Let
\[
 Y_R^s:=(V_R^{-s})^*.
\]

\begin{theorem}[Positive Sobolev regularity of the explicit complement]\label{thm:o3k-complement}
For the explicit vector
\[
 g_h=h-J_{R,S}u_h
\]
there exists a number
\[
 0<s_*=s_*(R,S,T_0,h)<1/2
\]
such that
\[
 \boxed{E_Ru_h\in H^{s_*}(\mathbb R),\qquad
 E_Sg_h\in H^{s_*}(\mathbb R).}
\tag{OJ.15}\label{eq:o3k-positive-sobolev}
\]
Consequently
\[
 \boxed{
 E_Sg_h\in\bigcap_{\alpha<\infty}\HG^\alpha(\mathbb R).
 }
\tag{OJ.16}\label{eq:o3k-all-log}
\]
In particular, if $m_h$ denotes the first nonzero integral-jet order of $g_h$, then
\[
 \boxed{E_Sg_h\in\HG^{m_h+3/2}(\mathbb R).}
\tag{OJ.17}\label{eq:o3k-log-gate}
\]
\end{theorem}

\begin{proof}
For $u\in V_R^s$ and $v\in V_R^{-s}$ define
\[
 \langle\mathcal A_Ru,v\rangle
 :=\mathfrak c_{\Gamma,R}[u,v]
 +\langle\Sigma_R^{[T_0]}u,v\rangle.
\tag{OJ.18}\label{eq:o3k-form-operator}
\]
The Gamma part is bounded by Fourier-space Cauchy--Schwarz:
\[
 |\mathfrak c_{\Gamma,R}[u,v]|
 \le \|u\|_{V_R^s}\|v\|_{V_R^{-s}}.
\]
By Lemma~\ref{lem:o3k-schur-window}, zero extension/restriction between the fixed nested
intervals and the factorization
\[
 \Sigma_{T_0}=H_{T_0}B_{T_0}H_{T_0}^*
\]
show, after shrinking $\delta>0$, that
$\Sigma_R^{[T_0]}$ is bounded on $H^s(-R,R)$ for $|s|\le\delta$.  Since $m\ge1$,
\[
 V_R^s\hookrightarrow H^s,
 \qquad V_R^{-s}\hookrightarrow H^{-s},
\]
so the $H^s$--$H^{-s}$ dual pairing gives
\[
 |\langle\Sigma_R^{[T_0]}u,v\rangle|
 \le C_s\|u\|_{V_R^s}\|v\|_{V_R^{-s}}.
\]
Thus
\[
 \mathcal A_R:V_R^s\longrightarrow Y_R^s
\]
is bounded throughout a symmetric interpolation interval around zero.

At $s=0$, $V_R^0$ is precisely the Gamma form space under zero extension, and
\[
 \langle\mathcal A_Ru,u\rangle
 =\mathfrak c_{\Gamma,R}[u]+\langle\Sigma_R^{[T_0]}u,u\rangle
 \ge \|u\|_{V_R^0}^2.
\]
Lax--Milgram therefore gives an isomorphism
\[
 \mathcal A_R:V_R^0\xrightarrow{\sim}Y_R^0.
\]
Apply the \v Sne\u\i berg local-stability theorem to the interpolation scales
\eqref{eq:o3k-interpolation} and their dual scales.  There exists $\varepsilon_1>0$ such
that
\[
 \mathcal A_R:V_R^s\xrightarrow{\sim}Y_R^s
 \qquad(|s|<\varepsilon_1).
\tag{OJ.19}\label{eq:o3k-form-invertible}
\]

By Proposition~\ref{prop:o3k-forcing}, $E_Rr_h\in H^{s_1}$ for some $s_1>0$.
Choose
\[
 0<s_*<\min\{s_1,\varepsilon_1\}.
\]
Then $r_h$ defines an element of $Y_R^{s_*}$ by the usual
$H^{s_*}$--$H^{-s_*}$ pairing.  Equation~\eqref{eq:o3k-form-invertible} therefore has a
unique solution $u_*\in V_R^{s_*}$.  Since $V_R^{s_*}\subset V_R^0$ and
$V_R^0\subset V_R^{-s_*}$, testing the $s_*$ equation against every vector in $V_R^0$
shows that $u_*$ solves the original weak equation~\eqref{eq:o3j-forcing-equation}.
Uniqueness at $s=0$ gives $u_*=u_h$.  Hence
\[
 E_Ru_h\in V_R^{s_*}\subset H^{s_*}(\mathbb R).
\]
Because
\[
 E_Sg_h=E_Sh-E_Ru_h
\]
and $E_Sh$ is Schwartz, the first assertion of~\eqref{eq:o3k-positive-sobolev} follows.
Finally, for every finite $\alpha$,
\[
 [\log(2+|\xi|)]^{2\alpha}
 \le C_{\alpha,s_*}\langle\xi\rangle^{2s_*},
\]
which proves~\eqref{eq:o3k-all-log} and hence~\eqref{eq:o3k-log-gate}.
\end{proof}

\begin{remark}[R12 firewall]\label{rem:o3k-firewall}
The exponent $s_*>0$ is a fixed-window quantity and may depend on $R,S,T_0$ and on the
chosen complement vector; no lower bound uniform in terminal horizons is asserted.
Theorem~\ref{thm:o3k-complement} closes the logarithmic regularity gate used by the O3
second-moment diagnostic.  It does not control the polar unitary gauges of
Remark~\ref{rem:polar-gauge}, does not prove $K_{R,S}^{T,U}\to I$, and does not prove or
disprove strong terminal transport.
\end{remark}

% P11 paper module: R13 polynomial second-moment witness and obstruction of the auxiliary Jensen product condition.

\subsection{Polynomial second-moment witness from the resolved complement gate}\label{subsec:o3l-polynomial-witness}

Keep $0<R<S<T_0$ fixed, let $g_h\in\mathcal C^-_{S,T_0}(R)$ be the explicit nonzero
odd complement above, and let $m_h$ be its first nonzero integral-jet order.  By
Theorem~\ref{thm:o3k-complement},
\[
 E_Sg_h\in\HG^{m_h+3/2}(\mathbb R),
\]
so Corollary~\ref{cor:log-threshold} applies to $g_h$.

Choose once and for all a smooth odd vector
\[
 f_0\in C_c^\infty((-R,R)),
 \qquad \beta_R^{(0)}(f_0)\ne0,
\]
and write
\[
 F_0:=J_{R,S}f_0.
\]
By the exact jet pullback, $F_0$ has first jet zero at level $S$.

\begin{lemma}[Mixed asymptotic with the rough complement]\label{lem:o3l-mixed-rough}
One has
\[
 \boxed{
 \sigma_U(J_{S,U}F_0,J_{S,U}g_h)
 =c_{m_h}\,\beta_R^{(0)}(f_0)
 \overline{\beta_S^{(m_h)}(g_h)}
 \frac{e^U}{U^{m_h+2}}(1+o(1)).
 }
\tag{O3L.1}\label{eq:o3l-mixed-rough}
\]
\end{lemma}

\begin{proof}
The rank-one decomposition in the proof of Theorem~\ref{thm:mixed-jet} is purely
Hilbert-space algebra and does not require smoothness.  For a fixed compactly supported
odd vector $f$, write
\[
 D_U(f,g)=\sigma_U(Jf,Jg)-
 \frac{\ell_U(f)\overline{\ell_U(g)}}{d_U}.
\]
Equation~\eqref{eq:mj-remainder-gram} gives the positive Gram representation and hence
\[
 |D_U(f,g)|^2\le D_U(f,f)D_U(g,g).
\tag{O3L.2}\label{eq:o3l-Dcs}
\]
For the smooth vector $F_0$, Theorem~\ref{thm:odd} gives the sharp diagonal asymptotic.
For $g_h$, Theorem~\ref{thm:o3k-complement} and
Corollary~\ref{cor:log-threshold} put $g_h$ within the scope of
\eqref{eq:rough-sharp-asymptotic}, so it has the corresponding sharp diagonal
asymptotic with first jet $m_h$.  The rough boundary expansion
\eqref{eq:rough-boundary-expansion} applies to both vectors, and
$d_U=2U+O(1)$ is independent of the source vector.  Therefore the rank-one diagonal
terms have the same leading terms as the corresponding total diagonal Schur forms.
It follows that
\[
 D_U(F_0,F_0)=o(e^U/U^2),
 \qquad
 D_U(g_h,g_h)=o(e^U/U^{2m_h+2}).
\]
By~\eqref{eq:o3l-Dcs},
\[
 D_U(F_0,g_h)=o(e^U/U^{m_h+2}).
\]
The rank-one term itself is obtained from the two boundary expansions and equals the
right-hand side of~\eqref{eq:o3l-mixed-rough} with relative error $O(U^{-1})$.  This
proves the lemma.
\end{proof}

\begin{lemma}[Crude relative-metric norm bound]\label{lem:o3l-relative-upper}
For $X=R,S$,
\[
 \boxed{
 \|A_{T_0,U}^{X,-}\|\le 1+\|H_U\|^2
 \ll \frac{e^U}{U}.
 }
\tag{O3L.3}\label{eq:o3l-relative-upper}
\]
\end{lemma}

\begin{proof}
For every nonzero odd $z\in\KXR{X}$, the Rayleigh quotient of the relative metric is
\[
 \frac{q_U^X(E_{X,U}z)}{q_{T_0}^X(E_{X,T_0}z)}.
\]
Exact Gamma compatibility, positivity of the Schur term, and
\eqref{eq:graph-norm-equiv} at level $U$ give
\[
 q_U^X(E_{X,U}z)
 \le (1+\|H_U\|^2)\,\mathfrak c_{\Gamma,X}[z],
 \qquad
 q_{T_0}^X(E_{X,T_0}z)\ge\mathfrak c_{\Gamma,X}[z].
\]
Hence the first inequality in~\eqref{eq:o3l-relative-upper} follows.

Furthermore
\[
 \|H_U\|
 \le2\sum_{p^k\le e^{2U}}\sqrt{\log p}\,p^{-3k/4}
 \le C\sum_{p\le e^{2U}}\sqrt{\log p}\,p^{-3/4}.
\]
The final prime sum is $O(e^{U/2}/\sqrt U)$ by the same Chebyshev/partial-summation
estimate used in~\eqref{eq:rough-kernel-modulus}.  Thus
$\|H_U\|^2\ll e^U/U$.
\end{proof}

\begin{theorem}[Polynomial second-moment witness]\label{thm:o3l-polynomial-witness}
For the actual P11 odd relative metrics at fixed $0<R<S<T_0$, there is a constant
$c>0$ such that, for all sufficiently large $U$,
\[
 \boxed{
 \nu_2(U)
 :=\frac{\|\Delta_2(U)\|}
 {\|A_{T_0,U}^{R,-}\|\,\|A_{T_0,U}^{S,-}\|}
 \ge c\,U^{-2m_h-2}.
 }
\tag{O3L.4}\label{eq:o3l-nu-lower}
\]
Consequently
\[
 \boxed{
 \|\Theta_{T_0,U}^-\|\ge \frac c8\,U^{-2m_h-2}.
 }
\tag{O3L.5}\label{eq:o3l-theta-polynomial}
\]
In particular, the auxiliary Jensen-product condition fails strongly:
\[
 \boxed{
 \chi_{T_0,U}^{R,-}\,\|\Theta_{T_0,U}^-\|\longrightarrow\infty.
 }
\tag{O3L.6}\label{eq:o3l-product-diverges}
\]
\end{theorem}

\begin{proof}
Set
\[
 A_R:=A_{T_0,U}^{R,-},\qquad
 A_S:=A_{T_0,U}^{S,-},\qquad
 W:=W_{R,S,-}^{[T_0]},
\]
and
\[
 \mathscr B=(I-WW^*)A_SW.
\]
Put
\[
 x:=G_{R,T_0}^{1/2}f_0,
 \qquad
 y:=G_{S,T_0}^{1/2}g_h.
\]
Because $g_h\in\ker(J_{R,S}^*G_{S,T_0})$, one has
$y\perp\Ran W$.  Therefore
\[
 \langle\mathscr Bx,y\rangle_{X,S}
 =\langle A_SWx,y\rangle_{X,S}.
\]
Using the definitions of $A_S$ and $W$ gives the exact identity
\[
 \boxed{
 \langle\mathscr Bx,y\rangle_{X,S}
 =\langle(G_{S,U}-G_{S,T_0})F_0,g_h\rangle_{X,S}.
 }
\tag{O3L.7}\label{eq:o3l-cross-exact}
\]
The subtraction of $G_{S,T_0}$ is allowed because the corresponding pairing is zero by
the complement condition.

The Gamma form is exactly invariant under zero extension, so the Gamma contributions in
\eqref{eq:o3l-cross-exact} cancel between $U$ and $T_0$.  Hence
\[
 \langle(G_{S,U}-G_{S,T_0})F_0,g_h\rangle_{X,S}
 =\sigma_U(J_{S,U}F_0,J_{S,U}g_h)
  -\sigma_{T_0}(J_{S,T_0}F_0,J_{S,T_0}g_h).
\]
The second term is fixed, while Lemma~\ref{lem:o3l-mixed-rough} has a nonzero leading
coefficient.  Thus
\[
 |\langle\mathscr Bx,y\rangle_{X,S}|
 \ge c_0\frac{e^U}{U^{m_h+2}}
\]
for all sufficiently large $U$.  Since $x$ and $y$ are fixed nonzero vectors,
\[
 \|\mathscr B\|
 \ge c_1\frac{e^U}{U^{m_h+2}}.
\tag{O3L.8}\label{eq:o3l-B-lower}
\]
By~\eqref{eq:jd-delta2},
\[
 \|\Delta_2\|=\|\mathscr B^*\mathscr B\|=\|\mathscr B\|^2
 \ge c_2\frac{e^{2U}}{U^{2m_h+4}}.
\tag{O3L.9}\label{eq:o3l-delta-lower}
\]
Lemma~\ref{lem:o3l-relative-upper} gives
\[
 \|A_R\|\,\|A_S\|\le C\frac{e^{2U}}{U^2},
\]
and division proves~\eqref{eq:o3l-nu-lower}.  The bound
\eqref{eq:o3l-theta-polynomial} follows from~\eqref{eq:jd-theta-lower}.

Finally Corollary~\ref{cor:conditioning} gives, for every $N>0$,
\[
 U^{-N}\chi_{T_0,U}^{R,-}\to\infty.
\]
Taking $N=2m_h+2$ and using~\eqref{eq:o3l-theta-polynomial} yields
\eqref{eq:o3l-product-diverges}.
\end{proof}

\begin{remark}[R13 firewall]\label{rem:o3l-firewall}
Theorem~\ref{thm:o3l-polynomial-witness} rules out the auxiliary product condition
$\chi_-\|\Theta_-\|\to0$ for this explicit fixed complement diagnostic.  By
Remark~\ref{rem:product} and the polar-gauge identity~\eqref{eq:jd-polar-gauge}, this does
\emph{not} imply failure of strong terminal transport.  It gives no control of the polar
unitaries $U_R,U_S$ and no conclusion about $K_{R,S}^{T,U}\to I$.
\end{remark}

% P11 paper module: R14 exact polar-gauge separation and no-promotion countermodel.

\subsection{Exact polar-gauge separation and algebraic non-promotion}\label{subsec:o3m-polar-gauge}

Keep the notation of Remark~\ref{rem:polar-gauge} at fixed
$0<R<S<T_0<U$ on the odd sectors:
\[
 W:=W_{R,S,-}^{[T_0]},
 \qquad
 Q:=A_S^{1/2}WA_R^{-1/2},
\]
\[
 X_R=(G_{R,U}^-)^{1/2}(G_{R,T_0}^-)^{-1/2},
 \qquad
 X_S=(G_{S,U}^-)^{1/2}(G_{S,T_0}^-)^{-1/2}.
\]
The polar decompositions are
\[
 X_R=U_RA_R^{1/2},
 \qquad
 X_S=U_SA_S^{1/2}.
\]

\begin{proposition}[Uniqueness of the polar gauges]\label{prop:o3m-polar-unique}
The unitary factors $U_R,U_S$ are uniquely determined by the corresponding metric
pairs.  Explicitly,
\[
 \boxed{
 U_R=X_RA_R^{-1/2},
 \qquad
 U_S=X_SA_S^{-1/2}.
 }
\tag{O3M.1}\label{eq:o3m-polar-unique}
\]
Thus the gauge issue is not a nonuniqueness of polar decomposition.
\end{proposition}

\begin{proof}
For a boundedly invertible operator $X$, the positive factor in its polar decomposition
is uniquely $|X|=(X^*X)^{1/2}$ and the unitary factor is uniquely
$X|X|^{-1}$.  Here $|X_R|=A_R^{1/2}$ and $|X_S|=A_S^{1/2}$.
\end{proof}

Let
\[
 K_{R,S}^{T_0,U}:=W^*W_{R,S,-}^{[U]}
\]
be the cross-terminal kernel.  Recall the modulus defects
\[
 \mathscr K:=I-W^*Q,
 \qquad
 \mathscr N:=(I-WW^*)Q,
\]
so that
\[
 Q=W(I-\mathscr K)+\mathscr N.
\]

\begin{proposition}[Exact gauge/modulus separation of the cross kernel]
\label{prop:o3m-cross-separation}
One has
\[
 \boxed{
 K_{R,S}^{T_0,U}
 =\Bigl[W^*U_SW(I-\mathscr K)+W^*U_S\mathscr N\Bigr]U_R^*.
 }
\tag{O3M.2}\label{eq:o3m-cross-separation}
\]
In particular, with
\[
 \Gamma_U:=W^*U_SWU_R^*,
\]
\[
 \boxed{
 \|K_{R,S}^{T_0,U}-\Gamma_U\|
 \le \|\mathscr K\|+\|\mathscr N\|
 \le 2\|Q-W\|.
 }
\tag{O3M.3}\label{eq:o3m-gauge-error}
\]
Consequently, even under the hypothetical stronger assumption $Q-W\to0$, convergence
of the actual cross kernel to the identity still requires the separate gauge condition
\[
 \Gamma_U\to I.
\]
\end{proposition}

\begin{proof}
The exact polar-gauge identity~\eqref{eq:jd-polar-gauge} gives
\[
 K_{R,S}^{T_0,U}=W^*U_SQU_R^*.
\]
Insert $Q=W(I-\mathscr K)+\mathscr N$ to obtain
\eqref{eq:o3m-cross-separation}.  Since $W$ is an isometry and $U_S,U_R$ are unitary,
\[
 \|K_{R,S}^{T_0,U}-\Gamma_U\|
 \le\|\mathscr K\|+\|\mathscr N\|.
\]
For every source vector $f$, the two summands in
$Q-W=-W\mathscr K+\mathscr N$ are orthogonal.  Hence
$\|\mathscr K\|\le\|Q-W\|$ and
$\|\mathscr N\|\le\|Q-W\|$, proving~\eqref{eq:o3m-gauge-error}.
\end{proof}

The next finite-dimensional model shows that no conclusion about the actual transport
can be extracted from the modulus data alone, even in the extreme case
$Q=W$ and $\Theta=0$ exactly.

\begin{proposition}[Exact algebraic countermodel to modulus-only promotion]
\label{prop:o3m-countermodel}
There exist finite-dimensional positive metric data satisfying the exact pullback,
relative-compression and polar-gauge identities for which
\[
 \boxed{Q=W=I,\qquad \Theta=0,}
\tag{O3M.4}\label{eq:o3m-perfect-modulus}
\]
but
\[
 \boxed{W^{[U]}\ne W,\qquad K^{T_0,U}\ne I.}
\tag{O3M.5}\label{eq:o3m-nontrivial-gauge}
\]
Moreover, with one fixed baseline system one can choose a sequence of future positive
metrics for which~\eqref{eq:o3m-perfect-modulus} holds at every step while the actual
future transports are not Cauchy.
\end{proposition}

\begin{proof}
Work on $\mathbb C^2$ and let
\[
 B=\begin{pmatrix}2&1\\1&2\end{pmatrix}>0,
 \qquad
 A=\begin{pmatrix}4&0\\0&1\end{pmatrix}>0.
\]
These matrices do not commute.  Put
\[
 G_{R,T_0}=I,
 \qquad
 G_{S,T_0}=B,
 \qquad
 J=B^{-1/2},
\]
and at the future horizon put
\[
 G_{R,U}=A,
 \qquad
 G_{S,U}=B^{1/2}AB^{1/2}.
\]
Then
\[
 J^*G_{S,T_0}J=G_{R,T_0},
 \qquad
 J^*G_{S,U}J=G_{R,U},
\]
so the exact terminal pullback identities hold.  The normalized base transport is
\[
 W=G_{S,T_0}^{1/2}JG_{R,T_0}^{-1/2}=I.
\]
The relative metrics are
\[
 A_R=G_{R,T_0}^{-1/2}G_{R,U}G_{R,T_0}^{-1/2}=A,
\]
\[
 A_S=G_{S,T_0}^{-1/2}G_{S,U}G_{S,T_0}^{-1/2}=A.
\]
Therefore
\[
 Q=A_S^{1/2}WA_R^{-1/2}=I,
\]
and the Jensen defect vanishes identically: $\Theta=0$.

On the other hand the actual future normalized transport is
\[
 V:=W^{[U]}
 =(B^{1/2}AB^{1/2})^{1/2}B^{-1/2}A^{-1/2}.
\tag{O3M.6}\label{eq:o3m-V}
\]
It is unitary by the pullback identity.  If $V=I$, then
\[
 (B^{1/2}AB^{1/2})^{1/2}=A^{1/2}B^{1/2}.
\]
The left side is positive selfadjoint, so the right side would be selfadjoint.  Hence
$A^{1/2}$ and $B^{1/2}$ would commute, equivalently $A$ and $B$ would commute, a
contradiction.  Thus $V\ne I$, proving~\eqref{eq:o3m-nontrivial-gauge}.

For the sequential statement keep the same fixed baseline data $B,J$ and alternate the
future source metric between $A$ and $I$.  Define the target future metric each time by
$B^{1/2}A_nB^{1/2}$.  At every step one still has $A_R=A_S=A_n$, hence
$Q=W=I$ and $\Theta=0$.  For $A_n=I$ the actual future transport is $I$, whereas for
$A_n=A$ it is the fixed nonidentity unitary $V$ in~\eqref{eq:o3m-V}.  The resulting
future transports therefore alternate between two distinct operators and are not
Cauchy.
\end{proof}

\begin{remark}[R14 firewall]\label{rem:o3m-firewall}
Proposition~\ref{prop:o3m-countermodel} is an algebraic non-promotion theorem: the exact
P11 identities by themselves do not imply terminal convergence from modulus/Jensen
data, even under the idealized condition $\Theta=0$.  It is not a counterexample to the
actual P11 metric family.  The actual asymptotic behavior of $U_R,U_S$ and
$K_{R,S}^{T_0,U}$ remains open and may still be constrained by additional arithmetic or
source structure not present in the abstract model.

In particular, the divergence
$\chi\|\Theta\|\to\infty$ proved in
Theorem~\ref{thm:o3l-polynomial-witness} only rules out the auxiliary sufficient
Jensen-product route.  It neither proves that $Q-W$ fails to tend to zero nor that the
actual cross-terminal kernel fails to tend to the identity.
\end{remark}

% P11 paper module: R15 exact gauge-intertwining criterion and square-root information gate.

\subsection{Gauge intertwining and the finite-jet square-root information gate}
\label{subsec:o3n-gauge-intertwining}

Keep $0<R<S<T_0<U$ fixed on the odd sectors and retain the notation
\[
 W:=W_{R,S,-}^{[T_0]},\qquad
 W_U:=W_{R,S,-}^{[U]},
\]
\[
 X_R=U_RA_R^{1/2},\qquad X_S=U_SA_S^{1/2},
\qquad
 Q=A_S^{1/2}WA_R^{-1/2},
\]
from Remarks~\ref{rem:polar-gauge} and Subsection~\ref{subsec:o3m-polar-gauge}.  Recall
\[
 W_U=U_SQU_R^*,
 \qquad
 \Gamma_U=W^*U_SWU_R^*.
\]

\begin{proposition}[Exact gauge-intertwining criterion]\label{prop:o3n-gauge-criterion}
Put
\[
 V_U:=U_SWU_R^*.
\]
Then $V_U$ and $W$ are isometries and
\[
 \Gamma_U=W^*V_U.
\]
For every source vector $f$,
\[
 \boxed{
 \|(V_U-W)f\|^2
 =2\|f\|^2-2\operatorname{Re}\langle\Gamma_Uf,f\rangle.
 }
\tag{O3N.1}\label{eq:o3n-gauge-rayleigh}
\]
Consequently
\[
 \boxed{
 \Gamma_U\xrightarrow[s]{}I
 \quad\Longleftrightarrow\quad
 V_U\xrightarrow[s]{}W.
 }
\tag{O3N.2}\label{eq:o3n-gauge-strong}
\]
In operator norm,
\[
 \boxed{
 \|\Gamma_U-I\|
 \le \|V_U-W\|
 \le \sqrt{2\|\Gamma_U-I\|},
 }
\tag{O3N.3}\label{eq:o3n-gauge-norm}
\]
so norm convergence of $\Gamma_U$ to $I$ is equivalent to norm convergence of the
transported gauge isometry $V_U$ to $W$.  Equivalently in operator norm,
\[
 \boxed{
 \Gamma_U\to I
 \quad\Longleftrightarrow\quad
 \|U_SW-WU_R\|\to0.
 }
\tag{O3N.4}\label{eq:o3n-gauge-intertwining}
\]
\end{proposition}

\begin{proof}
Since $U_R,U_S$ are unitary and $W$ is an isometry, so is $V_U$.  Therefore
\[
 (V_U-W)^*(V_U-W)=2I-W^*V_U-V_U^*W
 =2I-\Gamma_U-\Gamma_U^*,
\]
which gives~\eqref{eq:o3n-gauge-rayleigh}.  If $\Gamma_Uf\to f$ for a fixed $f$, the
right-hand side tends to zero, proving $V_Uf\to Wf$.  The converse follows from
$\Gamma_U-I=W^*(V_U-W)$.  This proves~\eqref{eq:o3n-gauge-strong}.

The same identity gives
\[
 \|V_U-W\|^2
 \le 2\|I-\Gamma_U\|,
\]
while $\|\Gamma_U-I\|\le\|V_U-W\|$, proving~\eqref{eq:o3n-gauge-norm}.  Finally
\[
 V_U-W=(U_SW-WU_R)U_R^*,
\]
so the two operator norms are equal.
\end{proof}

The preceding proposition isolates the genuine gauge question.  The unnormalised
square-root intertwining suggested by the metric products contains both the modulus and
gauge defects and is stronger than necessary.

\begin{proposition}[Exact square-root coherence decomposition]\label{prop:o3n-X-decomp}
One has
\[
 \boxed{
 X_SW-WX_R
 =\Bigl[U_S(Q-W)+(U_SW-WU_R)\Bigr]A_R^{1/2}.
 }
\tag{O3N.5}\label{eq:o3n-X-decomp}
\]
Equivalently, after the canonical right normalisation,
\[
 \boxed{
 \mathcal E_U
 :=(X_SW-WX_R)A_R^{-1/2}
 =U_SQ-WU_R.
 }
\tag{O3N.6}\label{eq:o3n-X-normalized}
\]
Moreover
\[
 \boxed{
 \mathcal E_UU_R^*=W_U-W,
 \qquad
 \|\mathcal E_U\|=\|W_U-W\|.
 }
\tag{O3N.7}\label{eq:o3n-X-transport}
\]
Thus the correctly normalised $X_S$--$X_R$ intertwining problem is exactly the original
norm terminal-transport problem; it is not an independent simplification.
\end{proposition}

\begin{proof}
From $A_S^{1/2}W=QA_R^{1/2}$ and the polar decompositions,
\[
 X_SW=U_SA_S^{1/2}W=U_SQA_R^{1/2},
 \qquad
 WX_R=WU_RA_R^{1/2}.
\]
Subtracting gives~\eqref{eq:o3n-X-decomp} and~\eqref{eq:o3n-X-normalized}.  Finally
$W_U=U_SQU_R^*$ gives~\eqref{eq:o3n-X-transport}.
\end{proof}

There is also an exact local interpretation of each individual polar gauge.

\begin{proposition}[Individual polar gauge equals metric noncommutativity]
\label{prop:o3n-gauge-commutator}
Let $X\in\{R,S\}$ and abbreviate
\[
 B:=G_{X,T_0}^-,\qquad C:=G_{X,U}^-.
\]
For the polar factor $U_X$ of
\[
 C^{1/2}B^{-1/2}=U_X
 \bigl(B^{-1/2}CB^{-1/2}\bigr)^{1/2},
\]
one has
\[
 \boxed{
 U_X=I
 \quad\Longleftrightarrow\quad
 [B,C]=0.
 }
\tag{O3N.8}\label{eq:o3n-gauge-commutator}
\]
Hence the word ``gauge'' records genuine noncommutativity between the fixed base metric
and the future metric; it is not a free phase choice.
\end{proposition}

\begin{proof}
If $B$ and $C$ commute, then $C^{1/2}B^{-1/2}$ is positive, so its polar unitary is
$I$.  Conversely, if $U_X=I$, then $C^{1/2}B^{-1/2}$ equals its positive polar modulus
and is therefore selfadjoint.  Hence
\[
 C^{1/2}B^{-1/2}=B^{-1/2}C^{1/2},
\]
so the positive square roots commute, equivalently $B$ and $C$ commute.
\end{proof}

The mixed-jet theorem determines all fixed-pair leading matrix elements, but its
leading matrix is asymptotically rank one.  The next elementary model shows why the
unquantified $o(1)$ remainders are insufficient for inverse-square-root control.

\begin{proposition}[Rank-one leading data do not determine inverse square roots]
\label{prop:o3n-rankone-insufficiency}
Let $\varepsilon\downarrow0$ and, for any fixed $a>2$, define
\[
 v_\varepsilon:=\binom{1}{\varepsilon},
 \qquad
 M_\varepsilon^{(a)}
 :=v_\varepsilon v_\varepsilon^*
   +\varepsilon^a e_2e_2^*
 =\begin{pmatrix}
 1&\varepsilon\\
 \varepsilon&\varepsilon^2+\varepsilon^a
 \end{pmatrix}>0.
\]
For every $a>2$ these matrices have exactly the same entrywise leading asymptotics
\[
 (M_\varepsilon^{(a)})_{11}=1,
 \qquad
 (M_\varepsilon^{(a)})_{12}=\varepsilon,
 \qquad
 (M_\varepsilon^{(a)})_{22}=\varepsilon^2(1+o(1)).
\tag{O3N.9}\label{eq:o3n-same-leading}
\]
However
\[
 \boxed{
 \lambda_{\min}(M_\varepsilon^{(a)})\sim\varepsilon^a,
 \qquad
 \|(M_\varepsilon^{(a)})^{-1/2}\|
 \sim\varepsilon^{-a/2}.
 }
\tag{O3N.10}\label{eq:o3n-inverse-root-scale}
\]
Thus identical rank-one fixed-pair leading data, even with relative $o(1)$ in every
entry, do not determine the polynomial scale of the inverse square root.
\end{proposition}

\begin{proof}
The determinant and trace are
\[
 \det M_\varepsilon^{(a)}=\varepsilon^a,
 \qquad
 \operatorname{tr}M_\varepsilon^{(a)}
 =1+\varepsilon^2+\varepsilon^a=1+o(1).
\]
Hence the larger eigenvalue tends to $1$, while the smaller eigenvalue, equal to the
determinant divided by the larger one, is asymptotic to $\varepsilon^a$.  The
inverse-square-root norm is the reciprocal square root of the smallest eigenvalue.
\end{proof}

The TC1 decomposition identifies the corresponding near-null direction inside the
actual fixed two-vector jet block.

\begin{proposition}[Exact TC1 near-null direction]\label{prop:o3n-nearnull}
Fix smooth odd vectors $f_0,f_m$ at one source level, with first nonzero integral-jet
orders $0$ and $m>0$, respectively.  For all sufficiently large $U$ one has
$\ell_U(f_0)\ne0$.  Define the $U$-dependent vector
\[
 \boxed{
 z_U:=f_m-\frac{\ell_U(f_m)}{\ell_U(f_0)}f_0.
 }
\tag{O3N.11}\label{eq:o3n-zU}
\]
Then the constant-mode rank-one form vanishes exactly on $z_U$,
\[
 \boxed{
 \rho_U(z_U,z_U)=0,
 \qquad
 \sigma_U(Jz_U,Jz_U)=D_U(z_U,z_U).
 }
\tag{O3N.12}\label{eq:o3n-nearnull-exact}
\]
Moreover
\[
 \boxed{
 D_U(z_U,z_U)
 =o\!\left(\frac{e^U}{U^{2m+2}}\right).
 }
\tag{O3N.13}\label{eq:o3n-nearnull-littleo}
\]
The present theory supplies no matching positive lower asymptotic for this quantity.
\end{proposition}

\begin{proof}
By~\eqref{eq:mj-ell},
\[
 \frac{\ell_U(f_m)}{\ell_U(f_0)}=O(U^{-m}),
\]
and $\ell_U(f_0)\ne0$ for all sufficiently large $U$.  The definition
\eqref{eq:o3n-zU} gives $\ell_U(z_U)=0$, hence
\[
 \rho_U(z_U,z_U)=\frac{|\ell_U(z_U)|^2}{d_U}=0.
\]
Since $\sigma_U=\rho_U+D_U$ on the fixed span of $f_0,f_m$, this proves
\eqref{eq:o3n-nearnull-exact}.

Expand $D_U(z_U,z_U)$ sesquilinearly in the fixed pair $f_0,f_m$.  Equations
\eqref{eq:mj-Ddiag}--\eqref{eq:mj-Dmixed}, together with
$\ell_U(f_m)/\ell_U(f_0)=O(U^{-m})$, show that each resulting term is
$o(e^U/U^{2m+2})$.  This proves~\eqref{eq:o3n-nearnull-littleo}.
\end{proof}

\begin{remark}[R15 information firewall]\label{rem:o3n-firewall}
On a two-vector jet-adapted block with first jet orders $0$ and $m>0$, the mixed-jet
asymptotic, after removal of its common scalar factor and a fixed diagonal
normalisation by the nonzero leading jet coefficients, has the rank-one scale
\[
 \begin{pmatrix}
 1&U^{-m}\\
 U^{-m}&U^{-2m}
 \end{pmatrix}
\]
up to entrywise relative $o(1)$ errors.  Proposition~\ref{prop:o3n-rankone-insufficiency}
shows that such information does not determine the smallest eigenvalue or the inverse
square root even to polynomial order.  Proposition~\ref{prop:o3n-nearnull} identifies
the exact P11 two-jet direction on which the leading rank-one term cancels: the missing
scale is the positive TC1 Gram remainder $D_U(z_U,z_U)$ (together with the fixed Gamma
floor in the full metric), not another copy of the already known leading boundary jet.

Therefore the present fixed-pair asymptotics cannot by themselves decide $U_R,U_S$,
$\Gamma_U$, or the true cross-terminal kernel.  This is an information-theoretic
obstruction to the current proof package, not a counterexample to the actual P11 metric
family.

Accordingly the genuine next step is exactly Open Problem~\ref{open:finite-jet-sqrt}:
one needs a quantitative finite-dimensional jet expansion deep enough to resolve the
near-null Gram directions and then compare the resulting square roots between the
$R$- and $S$-levels.  At the two-jet level, the sharpest immediate target is a
quantitative asymptotic or two-sided scale for $D_U(z_U,z_U)$.  No conclusion about
$\Gamma_U\to I$, $K_{R,S}^{T_0,U}\to I$, strong terminal transport, Object X, or RH is
drawn here.
\end{remark}

% P11 paper module: R16 absolute collapse of the TC1 near-null Schur remainder.

\subsection{Uniform collapse of the TC1 near-null Schur remainder}
\label{subsec:o3o-nearnull-collapse}

Keep the two fixed smooth odd source vectors $f_0,f_m\in C_c^\infty((-R,R))$
from Proposition~\ref{prop:o3n-nearnull}, with first nonzero integral-jet orders
$0$ and $m>0$.  Recall
\[
 z_U:=f_m-\frac{\ell_U(f_m)}{\ell_U(f_0)}f_0,
 \qquad \ell_U(z_U)=0.
\]
By~\eqref{eq:mj-ell},
\[
 \frac{\ell_U(f_m)}{\ell_U(f_0)}=O(U^{-m}),
\tag{O3O.1}\label{eq:o3o-coeff}
\]
so $z_U\to f_m$ in every fixed smooth seminorm.  In particular the family
$\{z_U\}$ has common compact support strictly inside $(-R,R)$ and uniformly bounded
$C^1$ and $L^2$ norms.

\begin{lemma}[Uniform R1 certificate on the fixed two-jet block]
\label{lem:o3o-uniform-certificate}
Let
\[
 h_U:=H_U^*J_{R,U}z_U,
 \qquad
 A_U:=I+R_U^*R_U.
\]
In the growing/remainder splitting of the proof of Theorem~\ref{thm:odd},
\[
 h_U=h_U^{\rm grow}+h_U^{\rm rem},
\]
one has uniformly
\[
 \boxed{\|h_U^{\rm rem}\|_2=O(1).}
\tag{O3O.2}\label{eq:o3o-rem-bounded}
\]
Writing
\[
 K_U:=\langle h_U^{\rm grow},\mathbf1_U\rangle,
 \qquad
 \mu_U:=\frac{K_U}{2U},
\]
one has
\[
 \boxed{|K_U|=O(\sqrt U),
 \qquad
 \|\mu_U\mathbf1_U\|_2^2=O(1).}
\tag{O3O.3}\label{eq:o3o-Kmu}
\]
Moreover the signed future-edge construction and the exact full-rest lift admit an
admissible certificate $\widehat Y_U^-$ and source remainder $Z_U^{\rm full}$ with
\[
 \boxed{
 \|\widehat Y_U^-\|^2=O(U^{-1}),
 \qquad
 \|Z_U^{\rm full}\|_2=O(1),
 }
\tag{O3O.4}\label{eq:o3o-certificate-cost}
\]
and
\[
 h_U=\widetilde R_U^*\widehat Y_U^-+Z_U^{\rm full}.
\tag{O3O.5}\label{eq:o3o-certificate-decomp}
\]
\end{lemma}

\begin{proof}
All estimates in Steps~2--6 of the proof of Theorem~\ref{thm:odd} depend on the
source through finitely many fixed support and smooth seminorm bounds.  By
\eqref{eq:o3o-coeff}, these are uniform on the bounded family $\{z_U\}$ inside the
fixed two-dimensional space $\operatorname{span}\{f_0,f_m\}$.  Hence
\eqref{eq:odd-rem-bounded} gives~\eqref{eq:o3o-rem-bounded} uniformly.

Since $\ell_U(z_U)=0$ exactly,
\[
 K_U
 =-\langle h_U^{\rm rem},\mathbf1_U\rangle,
\]
so Cauchy--Schwarz and~\eqref{eq:o3o-rem-bounded} yield
$|K_U|=O(\sqrt U)$ and therefore
\[
 \|\mu_U\mathbf1_U\|_2^2
 =\frac{|K_U|^2}{2U}=O(1).
\]

The kernel estimate~\eqref{eq:odd-kernel-bound} is uniform on $\{z_U\}$, hence
\eqref{eq:odd-kernel-weighted} gives the absolute bound
\[
 e^{-U}\int_0^Ue^{t/2}|k_U(t)|^2dt
 =O(U^{-1})+O(e^{-U/4}).
\]
For the mean subtraction,~\eqref{eq:o3o-Kmu} gives
\[
 e^{-U}\int_0^Ue^{t/2}\frac{|K_U|^2}{U^2}dt
 =O(e^{-U/2}/U).
\]
Thus the continuous signed certificate has squared norm $O(U^{-1})$.
The prime-cell mass/cost identity~\eqref{eq:prime-cell-cost} transfers the same
absolute bound to the discrete primitive certificate.  The source-representer
quadrature estimate~\eqref{eq:prime-cell-quadrature}, together with the exponentially
small future-cell mesh and the uniform derivative bound
\eqref{eq:odd-source-representer-derivative}, gives
\[
 \|Z_U^{\rm quad}\|_2=o(1).
\]
The constant part contributes only
$O(|K_U|e^{-2U/5})=o(1)$.  Finally
\eqref{eq:future-tail-small} and the $O(U^{-1})$ certificate cost give
\[
 \|Z_U^{\rm tail}\|_2=o(1).
\]
Using the same definitions as in~\eqref{eq:full-rest-lift}--
\eqref{eq:full-hub-cost},
\[
 Z_U^{\rm full}
 =\mu_U\mathbf1_U+h_U^{\rm rem}+Z_U^{\rm quad}+Z_U^{\rm tail},
\]
which proves~\eqref{eq:o3o-certificate-cost} and
\eqref{eq:o3o-certificate-decomp}.
\end{proof}

\begin{theorem}[Bounded TC1 near-null remainder]\label{thm:o3o-nearnull-bounded}
For the exact near-null vector $z_U$,
\[
 \boxed{
 D_U(z_U,z_U)
 =\sigma_U(J_{R,U}z_U,J_{R,U}z_U)
 =O(1).
 }
\tag{O3O.6}\label{eq:o3o-D-bounded}
\]
Consequently, for every fixed $N>0$,
\[
 \boxed{
 D_U(z_U,z_U)
 =o\!\left(\frac{e^U}{U^N}\right).
 }
\tag{O3O.7}\label{eq:o3o-superpoly-relative}
\]
In particular no constants $c>0$, $\alpha>0$ can satisfy
\[
 D_U(z_U,z_U)
 \ge c\frac{e^U}{U^{2m+2+\alpha}}
\]
for all sufficiently large $U$.
\end{theorem}

\begin{proof}
The exact dual formula~\eqref{eq:odd-dual} and
Lemma~\ref{lem:o3o-uniform-certificate} give
\[
 \sigma_U(J_{R,U}z_U,J_{R,U}z_U)
 \le
 \|Z_U^{\rm full}\|_2^2+\|\widehat Y_U^-\|^2
 =O(1).
\]
By construction $\ell_U(z_U)=0$, so the rank-one TC1 form
\[
 \rho_U(z_U,z_U)=\frac{|\ell_U(z_U)|^2}{d_U}
\]
vanishes exactly.  Since $D_U=\sigma_U-\rho_U$,
\eqref{eq:o3o-D-bounded} follows.  Equation
\eqref{eq:o3o-superpoly-relative} is immediate from
$U^Ne^{-U}\to0$.
\end{proof}

\begin{corollary}[Gamma floor on the exact near-null direction]
\label{cor:o3o-gamma-floor}
The full graph energy satisfies
\[
 \boxed{
 q_U^X(J_{R,U}z_U)\asymp1.
 }
\tag{O3O.8}\label{eq:o3o-full-floor}
\]
\end{corollary}

\begin{proof}
Terminal Gamma compatibility gives
\[
 q_U^X(J_{R,U}z_U)
 =\mathfrak c_{\Gamma,R}[z_U]+D_U(z_U,z_U).
\]
The second term is nonnegative and bounded by
Theorem~\ref{thm:o3o-nearnull-bounded}.  By~\eqref{eq:o3o-coeff},
$z_U\to f_m\ne0$, and the fixed Gamma form is positive, so
$\mathfrak c_{\Gamma,R}[z_U]$ stays bounded above and away from zero.
\end{proof}

\begin{corollary}[Exponential conditioning witness]
\label{cor:o3o-exp-conditioning}
At fixed $R$ and $T_0>R$,
\[
 \boxed{
 \kappa(A_{T_0,U}^{R,-})
 \ge c\frac{e^U}{U^2}
 }
\tag{O3O.9}\label{eq:o3o-exp-conditioning}
\]
for all sufficiently large $U$ and some $c>0$.
\end{corollary}

\begin{proof}
At the fixed baseline $T_0$,~\eqref{eq:o3o-coeff} gives
$z_U\to f_m$, so
\[
 q_{T_0}^X(J_{R,T_0}z_U)\asymp1.
\]
Together with~\eqref{eq:o3o-full-floor}, the relative Rayleigh quotient satisfies
\[
 \rho_{T_0,U}(z_U)=O(1).
\]
For the fixed first-jet-zero vector $f_0$, Theorem~\ref{thm:odd} gives
\[
 \rho_{T_0,U}(f_0)\asymp\frac{e^U}{U^2}.
\]
Using the Rayleigh characterization~\eqref{eq:conditioning-kappa},
\[
 \kappa(A_{T_0,U}^{R,-})
 \ge
 \frac{\rho_{T_0,U}(f_0)}{\rho_{T_0,U}(z_U)}
 \ge c\frac{e^U}{U^2}.
\]
\end{proof}

\begin{remark}[R16 firewall]\label{rem:o3o-firewall}
Theorem~\ref{thm:o3o-nearnull-bounded} determines the first absolute scale currently
provable after exact constant-mode cancellation: the Schur remainder is bounded, hence
smaller than $e^U/U^N$ for every fixed $N$.  It does \emph{not} determine a bounded-scale
limit.  A positive limit, decay to zero, or bounded oscillation are all still open.

The exponential conditioning bound~\eqref{eq:o3o-exp-conditioning} strengthens
Corollary~\ref{cor:conditioning}, but it supplies no asymptotic orientation of the
positive or inverse square roots and no control of the polar gauges.  In particular no
conclusion about
\[
 \Gamma_U\to I,
 \qquad
 K_{R,S}^{T_0,U}\to I,
\]
strong terminal transport, Object X, or RH follows.

The next genuine local target is the bounded core left by the certificate: determine
whether the residual contribution represented by $h_U^{\rm rem}$ (after the vanishing
future-edge certificate errors are removed) has a limit, a positive lower bound, or
further cancellation, and compare that bounded core compatibly at the $R$- and
$S$-levels.
\end{remark}

% P11 paper module: R17 vanishing of the bounded TC1 near-null Schur core.

\subsection{Vanishing of the bounded near-null Schur core}
\label{subsec:o3p-vanishing-core}

Keep the fixed smooth odd vectors $f_0,f_m\in C_c^\infty((-R,R))$ from
Proposition~\ref{prop:o3n-nearnull}, with first integral-jet orders $0$ and $m>0$, and
\[
 z_U=f_m-\frac{\ell_U(f_m)}{\ell_U(f_0)}f_0,
 \qquad \ell_U(z_U)=0.
\]
Recall from Theorem~\ref{thm:o3o-nearnull-bounded} that
\[
 D_U(z_U,z_U)=\sigma_U(J_{R,U}z_U,J_{R,U}z_U)=O(1).
\]
The bounded scale is not the final one: the residual source left by the first
future-edge certificate can itself be absorbed asymptotically.

\begin{proposition}[Global limit of the bounded hub remainder]
\label{prop:o3p-rem-limit}
Fix the cutoff $a_*$ used in the proof of Theorem~\ref{thm:odd}.  Define on the ambient
line
\[
 \mathcal H_{\rm rem,\infty}^*f
 :=
 \sum_{p:\,\frac12\log p<a_*}
 \sqrt{\log p}\,p^{-3/4}D_{\log p}^*E_Rf
 +
 \sum_p\sum_{k\ge2}
 \sqrt{\log p}\,p^{-3k/4}D_{k\log p}^*E_Rf.
\tag{O3P.1}\label{eq:o3p-global-rem}
\]
The second series converges absolutely in operator norm on both $L^1$ and $L^2$.
Moreover, after zero extension of the terminal remainder,
\[
 \boxed{
 E_Uh_U^{\rm rem}(z_U)
 \longrightarrow
 h_{\rm rem,\infty}:=\mathcal H_{\rm rem,\infty}^*f_m
 \quad\text{in }L^1(\mathbb R)\cap L^2(\mathbb R).
 }
\tag{O3P.2}\label{eq:o3p-rem-convergence}
\]
Finally,
\[
 \boxed{
 \int_{\mathbb R}h_{\rm rem,\infty}(x)\,dx=0.
 }
\tag{O3P.3}\label{eq:o3p-rem-meanzero}
\]
\end{proposition}

\begin{proof}
The primitive sum in~\eqref{eq:o3p-global-rem} is finite, while
\[
 \sum_p\sum_{k\ge2}\sqrt{\log p}\,p^{-3k/4}<\infty.
\]
Translations preserve $L^1$ and $L^2$ norms, so the higher-prime-power series converges
absolutely in operator norm on both spaces.  For every fixed $(p,k)$ the terminal cutoff
and terminal restriction disappear for all sufficiently large $U$ on the fixed compact
source support.  The omitted coefficient tail is uniformly small by absolute
summability, and $z_U\to f_m$ in the fixed smooth source space by
\eqref{eq:o3o-coeff}.  Dominated convergence over the coefficient family therefore gives
\eqref{eq:o3p-rem-convergence}.  Each complete full-line translation difference has
zero integral, and the series is absolutely $L^1$-convergent, proving
\eqref{eq:o3p-rem-meanzero}.
\end{proof}

Use the R16 notation
\[
 K_U=\langle h_U^{\rm grow},\mathbf1_U\rangle,
 \qquad
 \mu_U=\frac{K_U}{2U}.
\]
Since $\ell_U(z_U)=0$,
\[
 K_U=-\langle h_U^{\rm rem}(z_U),\mathbf1_U\rangle.
\]
Proposition~\ref{prop:o3p-rem-limit} therefore gives
\[
 \boxed{K_U\to0,
 \qquad
 \|\mu_U\mathbf1_U\|_2^2=\frac{|K_U|^2}{2U}\to0.}
\tag{O3P.4}\label{eq:o3p-mu-vanish}
\]
In particular, for
\[
 r_U:=\mu_U\mathbf1_U+h_U^{\rm rem}(z_U),
\]
one has after zero extension
\[
 \boxed{E_Ur_U\to h_{\rm rem,\infty}\quad\text{in }L^2(\mathbb R).}
\tag{O3P.5}\label{eq:o3p-r-limit}
\]

\begin{lemma}[Fixed mean-zero profiles are asymptotically absorbed]
\label{lem:o3p-fixed-absorption}
Let $g\in C_c^\infty(\mathbb R)$ be even with $\int g=0$.  Regard $g$ as a terminal
vector in $L^2(-U,U)$ for all sufficiently large $U$.  Then there exist
$Y_{U,g}\in\mathscr Z_U$ and $E_{U,g}\in L^2(-U,U)$ such that
\[
 \boxed{
 g=\widetilde R_U^*Y_{U,g}+E_{U,g},
 \qquad
 \|Y_{U,g}\|_{\mathscr Z_U}\to0,
 \qquad
 \|E_{U,g}\|_2\to0.
 }
\tag{O3P.6}\label{eq:o3p-fixed-cert}
\]
Consequently
\[
 \boxed{
 \langle g,(I+R_U^*R_U)^{-1}g\rangle\to0.
 }
\tag{O3P.7}\label{eq:o3p-fixed-energy}
\]
\end{lemma}

\begin{proof}
Repeat the signed future-edge construction in Steps~4--6 of the proof of
Theorem~\ref{thm:odd}.  For an even terminal test vector $e$, set
$b_U(t)=e(U-t)$ and
\[
 k_g^{(U)}(t):=2g(U-t),\qquad 0<t<U.
\]
The zero integral of $g$ gives $\int_0^Uk_g^{(U)}(t)dt=0$, so the exact anchor identity
\eqref{eq:odd-anchor-identity} applies.  The continuous primitive certificate cost is
bounded by
\[
 e^{-U}\int_0^Ue^{t/2}|k_g^{(U)}(t)|^2dt
 \le e^{-U/2}\|k_g^{(U)}\|_2^2
 \ll_g e^{-U/2}.
\tag{O3P.8}\label{eq:o3p-fixed-cont-cost}
\]
Because $g$ is fixed smooth and compactly supported, the source-representer derivative
is uniformly bounded.  All required future cells satisfy
$r\le(U+O(1))/2$, hence the $\theta=3/5$ cells of
\eqref{eq:prime-cell-mass} have maximal width $O(e^{-2U/5})$.
The mass-normalized quadrature~\eqref{eq:prime-cell-cost}--
\eqref{eq:prime-cell-quadrature} therefore has vanishing source error and vanishing
certificate cost.  The exact $a=0$ full-rest lift has vanishing higher-prime-power tail
by~\eqref{eq:future-tail-small}.  This gives~\eqref{eq:o3p-fixed-cert}; inserting it into
the dual formula~\eqref{eq:odd-dual} gives~\eqref{eq:o3p-fixed-energy}.
\end{proof}

\begin{theorem}[Vanishing TC1 near-null Schur core]
\label{thm:o3p-nearnull-vanish}
For the exact near-null vector $z_U$,
\[
 \boxed{
 D_U(z_U,z_U)
 =\sigma_U(J_{R,U}z_U,J_{R,U}z_U)
 \longrightarrow0.
 }
\tag{O3P.9}\label{eq:o3p-D-vanish}
\]
\end{theorem}

\begin{proof}
The first R16 certificate gives
\[
 h_U
 =\widetilde R_U^*\widehat Y_U^{(1)}
 +r_U+Z_U^{\rm quad}+Z_U^{\rm tail},
\tag{O3P.10}\label{eq:o3p-first-decomp}
\]
where
\[
 \|\widehat Y_U^{(1)}\|^2=O(U^{-1}),
 \qquad
 \|Z_U^{\rm quad}\|_2+\|Z_U^{\rm tail}\|_2=o(1).
\tag{O3P.11}\label{eq:o3p-first-cost}
\]
Fix $\varepsilon>0$.  By the absolute $L^2$ convergence in
Proposition~\ref{prop:o3p-rem-limit}, choose a finite partial sum
$g_\varepsilon$ of~\eqref{eq:o3p-global-rem} such that
\[
 \|h_{\rm rem,\infty}-g_\varepsilon\|_2<\varepsilon.
\]
It is a finite sum of complete translation differences, hence is even, smooth,
compactly supported and has zero integral.  By~\eqref{eq:o3p-r-limit}, for all
sufficiently large $U$,
\[
 \|r_U-g_\varepsilon\|_2<2\varepsilon.
\tag{O3P.12}\label{eq:o3p-r-approx}
\]
Apply Lemma~\ref{lem:o3p-fixed-absorption} to this fixed $g_\varepsilon$:
\[
 g_\varepsilon=\widetilde R_U^*Y_{U,\varepsilon}+E_{U,\varepsilon},
 \qquad
 \|Y_{U,\varepsilon}\|\to0,
 \qquad
 \|E_{U,\varepsilon}\|_2\to0.
\]
Insert this into~\eqref{eq:o3p-first-decomp} and use the dual formula
\eqref{eq:odd-dual}.  Taking first $U\to\infty$ with $\varepsilon$ fixed and then
$\varepsilon\downarrow0$ gives
\[
 \sigma_U(J_{R,U}z_U,J_{R,U}z_U)\to0.
\]
Since $\ell_U(z_U)=0$, the rank-one form $\rho_U(z_U,z_U)$ vanishes identically, so
$D_U=\sigma_U$ on this vector and~\eqref{eq:o3p-D-vanish} follows.
\end{proof}

\begin{corollary}[Pure Gamma limit on the exact near-null direction]
\label{cor:o3p-gamma-limit}
One has
\[
 \boxed{
 q_U^X(J_{R,U}z_U)
 \longrightarrow
 \mathfrak c_{\Gamma,R}[f_m]>0.
 }
\tag{O3P.13}\label{eq:o3p-gamma-limit}
\]
At fixed baseline $T_0>R$,
\[
 \boxed{
 \rho_{T_0,U}(z_U)
 \longrightarrow
 \frac{\mathfrak c_{\Gamma,R}[f_m]}
 {q_{T_0}^X(J_{R,T_0}f_m)}\in(0,\infty).
 }
\tag{O3P.14}\label{eq:o3p-rayleigh-limit}
\]
\end{corollary}

\begin{proof}
Terminal Gamma compatibility gives
\[
 q_U^X(J_{R,U}z_U)
 =\mathfrak c_{\Gamma,R}[z_U]+D_U(z_U,z_U).
\]
Now $z_U\to f_m$ in the fixed finite-dimensional smooth source space and
Theorem~\ref{thm:o3p-nearnull-vanish} gives $D_U(z_U,z_U)\to0$, proving
\eqref{eq:o3p-gamma-limit}.  The fixed-baseline denominator converges to
$q_{T_0}^X(J_{R,T_0}f_m)$, proving~\eqref{eq:o3p-rayleigh-limit}.
\end{proof}

\begin{proposition}[Source-compatible scalar cores carry no gauge difference]
\label{prop:o3p-source-compatible}
Let $0<R<S$ and choose the fixed pair $f_0,f_m$ in $C_c^\infty((-R,R))$.  Regard the
same vectors at level $S$ by zero extension.  Then for every $U>S$ the corresponding
exact near-null combinations satisfy
\[
 \boxed{z_U^S=J_{R,S}z_U^R.}
\tag{O3P.15}\label{eq:o3p-z-compatible}
\]
Moreover the scalar terminal TC1 remainders agree identically:
\[
 \boxed{
 D_U^{(S)}(z_U^S,z_U^S)
 =D_U^{(R)}(z_U^R,z_U^R).
 }
\tag{O3P.16}\label{eq:o3p-D-compatible}
\]
Thus a difference of these source-compatible scalar bounded cores cannot encode the
$R$--$S$ polar-gauge commutator.
\end{proposition}

\begin{proof}
The cocycle gives $J_{S,U}J_{R,S}=J_{R,U}$, so the terminal functional
$\ell_U(f)=\langle H_U^*J_{\cdot,U}f,\mathbf1_U\rangle$ has the same value for the same
terminal vector independent of whether it is regarded as originating at level $R$ or
$S$.  Hence the coefficient defining $z_U$ is identical, proving
\eqref{eq:o3p-z-compatible}.  The form $D_U$ is defined entirely at terminal level $U$
from $H_U$, $R_U$, $A_U$ and $\mathbf1_U$, so it has the same value on the same terminal
vector, proving~\eqref{eq:o3p-D-compatible}.
\end{proof}

\begin{remark}[R17 firewall]\label{rem:o3p-firewall}
Theorem~\ref{thm:o3p-nearnull-vanish} resolves the scalar bounded-core alternatives left
open in R16: the TC1 Schur remainder tends to zero, while the full graph energy tends to
the fixed positive Gamma floor.  Proposition~\ref{prop:o3p-source-compatible} also shows
that comparing this scalar core between canonically nested source levels cannot measure
the polar gauge.

The remaining terminal problem is therefore genuinely operator-valued.  It requires
control of off-diagonal finite-jet Gram data, the fixed Gamma matrix, and the orientation
of positive/inverse square roots under $R\to S$.  No conclusion about
\[
 \Gamma_U\to I,
 \qquad
 K_{R,S}^{T_0,U}\to I,
 \qquad
 W_{R,S,-}^{[U]}\text{ strong Cauchy},
\]
Object X, Seal, or RH follows here.
\end{remark}

% P11 paper module: R18 off-diagonal near-null block and finite source-block firewall.

\subsection{Off-diagonal near-null suppression and the finite source-block firewall}
\label{subsec:o3q-offdiagonal-nearnull}

Keep the fixed smooth odd pair $f_0,f_m\in C_c^\infty((-R,R))$ from
Proposition~\ref{prop:o3n-nearnull}, with first integral-jet orders $0$ and $m>0$, and
\[
 z_U=f_m-\frac{\ell_U(f_m)}{\ell_U(f_0)}f_0,
 \qquad \ell_U(z_U)=0.
\]
Recall from Theorem~\ref{thm:o3p-nearnull-vanish} that
\[
 D_U(z_U,z_U)=\sigma_U(J_{R,U}z_U,J_{R,U}z_U)\to0.
\]
The mixed entry on the pair $(f_0,z_U)$ is then forced below the natural square-root
boundary scale.

\begin{theorem}[Off-diagonal TC1 suppression]\label{thm:o3q-mixed-suppression}
One has
\[
 \boxed{
 \rho_U(f_0,z_U)=0,
 \qquad
 \sigma_U(J_{R,U}f_0,J_{R,U}z_U)
 =o\!\left(\frac{e^{U/2}}{U}\right).
 }
\tag{O3Q.1}\label{eq:o3q-mixed-suppression}
\]
\end{theorem}

\begin{proof}
Since $\ell_U(z_U)=0$,
\[
 \rho_U(f_0,z_U)
 =\frac{\ell_U(f_0)\overline{\ell_U(z_U)}}{d_U}=0,
\]
so
\[
 \sigma_U(Jf_0,Jz_U)=D_U(f_0,z_U).
\]
The positive Gram representation~\eqref{eq:mj-remainder-gram} gives
\[
 |D_U(f_0,z_U)|^2
 \le D_U(f_0,f_0)D_U(z_U,z_U).
\]
For the fixed first-jet-zero vector $f_0$, equation~\eqref{eq:mj-Ddiag} gives
\[
 D_U(f_0,f_0)=o(e^U/U^2),
\]
while Theorem~\ref{thm:o3p-nearnull-vanish} gives
$D_U(z_U,z_U)\to0$.  Hence
\[
 |D_U(f_0,z_U)|^2=o(e^U/U^2),
\]
which proves~\eqref{eq:o3q-mixed-suppression}.
\end{proof}

The full graph form has a stronger geometric consequence because R17 leaves a fixed
positive Gamma floor on $z_U$.

\begin{corollary}[Asymptotic full-Gram orthogonality]\label{cor:o3q-full-angle}
Put
\[
 L_U:=\frac{e^U}{U^2}.
\]
There is $a_0>0$ such that
\[
 q_U^X(J_{R,U}f_0)=a_0L_U(1+o(1)),
\tag{O3Q.2}\label{eq:o3q-f0-diag}
\]
while
\[
 q_U^X(J_{R,U}z_U)
 \longrightarrow
 \gamma_m:=\mathfrak c_{\Gamma,R}[f_m]>0.
\tag{O3Q.3}\label{eq:o3q-z-diag}
\]
Moreover
\[
 q_U^X(J_{R,U}f_0,J_{R,U}z_U)=o(L_U^{1/2}),
\tag{O3Q.4}\label{eq:o3q-full-mixed}
\]
and therefore
\[
 \boxed{
 \frac{q_U^X(J_{R,U}f_0,J_{R,U}z_U)}
 {q_U^X(J_{R,U}f_0)^{1/2}
  q_U^X(J_{R,U}z_U)^{1/2}}
 \longrightarrow0.
 }
\tag{O3Q.5}\label{eq:o3q-full-angle}
\]
Equivalently, for the coordinate Gram matrix
\[
 M_U:=
 \begin{pmatrix}
 q_U^X(Jf_0) & q_U^X(Jf_0,Jz_U)\\
 q_U^X(Jz_U,Jf_0) & q_U^X(Jz_U)
 \end{pmatrix}
\]
and
\[
 S_U:=\operatorname{diag}
 \bigl(q_U^X(Jf_0)^{-1/2},q_U^X(Jz_U)^{-1/2}\bigr),
\]
one has
\[
 \boxed{S_UM_US_U\to I_2,}
\tag{O3Q.6}\label{eq:o3q-correlation-matrix}
\]
and
\[
 \boxed{\det M_U\sim a_0\gamma_mL_U.}
\tag{O3Q.7}\label{eq:o3q-det}
\]
\end{corollary}

\begin{proof}
Theorem~\ref{thm:odd} with first jet $0$ gives the Schur asymptotic in
\eqref{eq:o3q-f0-diag}; the fixed Gamma term is $O(1)$ and hence negligible on the
scale $L_U$.  Equation~\eqref{eq:o3q-z-diag} is
Corollary~\ref{cor:o3p-gamma-limit}.

For the mixed form, terminal Gamma compatibility gives
\[
 q_U^X(Jf_0,Jz_U)
 =\mathfrak c_{\Gamma,R}[f_0,z_U]
  +\sigma_U(Jf_0,Jz_U).
\]
Since $z_U\to f_m$ in the fixed smooth source space,
$\mathfrak c_{\Gamma,R}[f_0,z_U]=O(1)$.  Theorem~\ref{thm:o3q-mixed-suppression}
therefore proves~\eqref{eq:o3q-full-mixed}.  Dividing by the two diagonal asymptotics
gives~\eqref{eq:o3q-full-angle}, and the matrix statements follow immediately.
\end{proof}

The preceding two-vector statement does not create a source-level gauge diagnostic.
In fact the cocycle trivializes every finite block made only from canonically nested
source vectors.

\begin{proposition}[Complete finite source-block cocycle invariance]
\label{prop:o3q-finite-block-invariance}
Let $0<R<S<T$ and choose any finite family
$e_1,\dots,e_n\in\KXR{R}$.  Put
\[
 e_i^S:=J_{R,S}e_i.
\]
Then for every terminal $T>S$ and every $i,j$,
\[
 \boxed{
 \langle G_{S,T}e_i^S,e_j^S\rangle_{X,S}
 =
 \langle G_{R,T}e_i,e_j\rangle_{X,R}.
 }
\tag{O3Q.8}\label{eq:o3q-block-invariance}
\]
Consequently, for fixed $T_0$ and future terminal $U$, the complete finite Gram pair
\[
 \left(
 [\langle G_{R,T_0}e_i,e_j\rangle],
 [\langle G_{R,U}e_i,e_j\rangle]
 \right)
\]
is identical, under the canonical coordinate identification, to the corresponding
pair after transport to level $S$.

The same assertion holds pointwise for $U$-dependent finite families, provided their
$S$-level representatives are the canonical images of the $R$-level vectors.
\end{proposition}

\begin{proof}
By the graph cocycle,
\[
 J_{S,T}e_i^S
 =J_{S,T}J_{R,S}e_i
 =J_{R,T}e_i.
\]
Since $G_{X,T}=J_{X,T}^*J_{X,T}$ with the graph-Hilbert adjoint,
\[
 \langle G_{S,T}e_i^S,e_j^S\rangle_{X,S}
 =\langle J_{S,T}e_i^S,J_{S,T}e_j^S\rangle_{X,T}
 =\langle J_{R,T}e_i,J_{R,T}e_j\rangle_{X,T},
\]
which is~\eqref{eq:o3q-block-invariance}.  Applying it at $T_0$ and at $U$ proves the
Gram-pair statement.  The proof is pointwise in the chosen vectors and therefore also
applies to a $U$-dependent family.
\end{proof}

\begin{corollary}[The $(f_0,z_U)$ mixed scalar is not an $R/S$ gauge diagnostic]
\label{cor:o3q-mixed-no-gauge}
Let $0<R<S$ and regard the pair $(f_0,z_U)$ at level $S$ by canonical zero extension.
Then its complete terminal $(T_0,U)$ Gram pair, including every off-diagonal entry, is
identical at the two source levels.  In particular no difference of
\[
 q_U^X(Jf_0,Jz_U),
 \qquad
 \sigma_U(Jf_0,Jz_U),
\]
or of any finite generalized Gram quantity built solely from the canonically nested
pair can measure the $R$--$S$ polar-gauge defect.
\end{corollary}

\begin{proof}
Apply Proposition~\ref{prop:o3q-finite-block-invariance} to the two-vector family.  The
Schur statement follows as well because the Gamma form is terminal-compatible, so
subtracting the identical Gamma matrix from the identical full terminal Gram matrix
leaves the identical Schur matrix.
\end{proof}

\begin{remark}[R18 firewall: internal blocks versus full square roots]
\label{rem:o3q-firewall}
Corollary~\ref{cor:o3q-full-angle} resolves the normalized orientation of the
\emph{internal two-vector future Gram block}: after exact boundary-mode cancellation the
large $f_0$ direction and the Gamma-supported $z_U$ direction are asymptotically
orthogonal.  Proposition~\ref{prop:o3q-finite-block-invariance} shows that this complete
finite form pair is already source-compatible under $R\to S$.

This does not determine the full polar gauges.  Compression and positive functional
calculus do not commute in general:
\[
 P_EA^{1/2}P_E\ne(P_EAP_E)^{1/2}.
\]
The canonical-inclusion model of Subsection~\ref{subsec:o3m-polar-gauge} already shows
that perfect internal modulus data need not determine the true future transport.
Therefore the remaining gauge information must involve the coupling of the nested
source image with its complement before the full square root/polar factor is taken.

A genuine next target is consequently an off-block square-root quantity in the
$T_0$-orthogonal splitting
\[
 \Ran W\oplus(\Ran W)^\perp,
\]
for example the relation between the already nontrivial block
\[
 (I-WW^*)A_SW
\]
and its square-root analogue
\[
 (I-WW^*)A_S^{1/2}W,
\]
or directly the polar leakage $(I-WW^*)U_SW$.
No conclusion about $\Gamma_U\to I$, $K_{R,S}^{T_0,U}\to I$, strong terminal
transport, Object X, Seal, or RH is drawn here.
\end{remark}

% P11 paper module: R19 square-root off-block witness and modulus/polar separation.

\subsection{Square-root off-block witness and modulus/polar separation}
\label{subsec:o3r-sqrt-offblock}

Keep the odd-sector notation at fixed $0<R<S<T_0<U$,
\[
 A_R:=A_{T_0,U}^{R,-},\qquad
 A_S:=A_{T_0,U}^{S,-},\qquad
 W:=W_{R,S,-}^{[T_0]},\qquad P:=WW^*.
\]
Put
\[
 \mathscr B_U:=(I-P)A_SW,
 \qquad
 \mathscr C_U:=(I-P)A_S^{1/2}W.
\tag{O3R.1}\label{eq:o3r-BC}
\]
The first operator is the second-moment off-block from
\eqref{eq:jd-B}; the second is its square-root analogue.

\begin{proposition}[Exact square-root block factorization]
\label{prop:o3r-block-factorization}
Let
\[
 D_U:=W^*A_S^{1/2}W,
 \qquad
 E_U:=(I-P)A_S^{1/2}(I-P)|_{\Ran W^\perp}.
\]
Then
\[
 \boxed{\mathscr B_U=\mathscr C_UD_U+E_U\mathscr C_U.}
\tag{O3R.2}\label{eq:o3r-factorization}
\]
Consequently
\[
 \boxed{
 \|\mathscr C_U\|
 \ge
 \frac{\|\mathscr B_U\|}
 {\sqrt{\|A_R\|}+\sqrt{\|A_S\|}}.
 }
\tag{O3R.3}\label{eq:o3r-C-lower-general}
\]
\end{proposition}

\begin{proof}
This is the lower-bound direction of the block identity already used in the proof of
Proposition~\ref{prop:second}.  Writing $S_U=A_S^{1/2}$, one has
$PS_UW=WD_U$ and therefore
\[
 (I-P)S_U^2W
 =(I-P)S_UPS_UW+(I-P)S_U(I-P)S_UW
 =\mathscr C_UD_U+E_U\mathscr C_U.
\]
This proves~\eqref{eq:o3r-factorization}.  Operator Jensen gives
$0\le D_U\le A_R^{1/2}$, so
$\|D_U\|\le\sqrt{\|A_R\|}$, while
$\|E_U\|\le\sqrt{\|A_S\|}$.  Hence
\[
 \|\mathscr B_U\|
 \le(\|D_U\|+\|E_U\|)\|\mathscr C_U\|,
\]
which yields~\eqref{eq:o3r-C-lower-general}.
\end{proof}

\begin{theorem}[P11 square-root off-block witness]
\label{thm:o3r-sqrt-witness}
Let $m_h$ be the first nonzero integral-jet order of the fixed complement vector
$g_h$ used in Theorem~\ref{thm:o3l-polynomial-witness}.  Then there exists $c>0$ such
that for all sufficiently large $U$,
\[
 \boxed{
 \|\mathscr C_U\|
 \ge c\frac{e^{U/2}}{U^{m_h+3/2}}.
 }
\tag{O3R.4}\label{eq:o3r-C-lower}
\]
In particular the square-root off-block diverges in norm.
\end{theorem}

\begin{proof}
Equation~\eqref{eq:o3l-B-lower} gives
\[
 \|\mathscr B_U\|\ge c_0\frac{e^U}{U^{m_h+2}}.
\]
By Lemma~\ref{lem:o3l-relative-upper},
\[
 \|A_R\|,\|A_S\|\ll\frac{e^U}{U}.
\]
Insert these bounds in~\eqref{eq:o3r-C-lower-general}.
\end{proof}

Recall the modulus isometry and its off-range defect
\[
 Q=A_S^{1/2}WA_R^{-1/2},
 \qquad
 \mathscr N_U:=(I-P)Q.
\]

\begin{corollary}[Polynomial modulus-leakage lower bound]
\label{cor:o3r-modulus-leakage}
One has the exact identity
\[
 \boxed{\mathscr C_U=\mathscr N_UA_R^{1/2}}
\tag{O3R.5}\label{eq:o3r-CN}
\]
and therefore
\[
 \boxed{
 \|\mathscr N_U\|\ge cU^{-m_h-1}.
 }
\tag{O3R.6}\label{eq:o3r-N-lower}
\]
Consequently
\[
 \boxed{
 \|Q-W\|\ge cU^{-m_h-1}.
 }
\tag{O3R.7}\label{eq:o3r-QW-lower}
\]
Thus the auxiliary modulus isometry cannot approach $W$ faster than this polynomial
scale in operator norm.
\end{corollary}

\begin{proof}
Equation~\eqref{eq:o3r-CN} follows directly from
$\mathscr N_U=(I-P)A_S^{1/2}WA_R^{-1/2}$.  Hence
\[
 \|\mathscr C_U\|
 \le\|\mathscr N_U\|\sqrt{\|A_R\|}.
\]
Combine Theorem~\ref{thm:o3r-sqrt-witness} with
$\sqrt{\|A_R\|}\ll e^{U/2}U^{-1/2}$ to obtain
\eqref{eq:o3r-N-lower}.  Finally
\[
 Q-W=-W\mathscr K+\mathscr N_U
\]
is an orthogonal range decomposition, so
$\|Q-W\|\ge\|\mathscr N_U\|$.
\end{proof}

The preceding result is genuinely a square-root/modulus statement.  It must not be
promoted to the polar unitary factors.

\begin{proposition}[Square-root leakage does not algebraically force polar leakage]
\label{prop:o3r-polar-countermodel}
There exists a canonical inclusion model satisfying the exact relative-metric
compression identity in which $\mathscr B_U$ and $\mathscr C_U$ are nonzero (indeed can
be made arbitrarily large), while
\[
 \boxed{(I-P)U_SW=0}
\tag{O3R.8}\label{eq:o3r-polar-zero}
\]
identically.
\end{proposition}

\begin{proof}
Take source space $\mathbb C$, target space $\mathbb C^2$, and
\[
 Wz=(z,0),\qquad P=WW^*.
\]
Let the base metrics be
$G_{R,T_0}=1$ and $G_{S,T_0}=I_2$.  Put
\[
 A_0=\begin{pmatrix}2&1\\1&2\end{pmatrix}>0,
 \qquad
 G_{S,U}=t^2A_0,
 \qquad
 G_{R,U}=2t^2
\]
with $t>0$.  Then the relative metrics are
$A_S=t^2A_0$ and $A_R=2t^2$, and
\[
 W^*A_SW=A_R
\]
exactly.  Moreover
\[
 X_S=G_{S,U}^{1/2}G_{S,T_0}^{-1/2}=tA_0^{1/2},
 \qquad
 X_R=\sqrt2\,t,
\]
are positive.  Their polar factors are therefore
$U_S=I_2$ and $U_R=1$, proving~\eqref{eq:o3r-polar-zero}.

On the other hand,
\[
 \mathscr B_U=t^2(I-P)A_0W\ne0,
 \qquad
 \mathscr C_U=t(I-P)A_0^{1/2}W\ne0,
\]
because $e_1$ is not an eigenvector of $A_0$.  Thus the two norms grow like $t^2$ and
$t$, respectively, while the polar leakage remains exactly zero.
\end{proof}

\begin{remark}[R19 firewall]\label{rem:o3r-firewall}
Theorem~\ref{thm:o3r-sqrt-witness} is the first direct P11 lower witness obtained after
passing through the positive square root of the future relative metric.  Corollary
\ref{cor:o3r-modulus-leakage} strengthens the earlier Jensen diagnostics by giving a
direct quantitative lower bound on the actual modulus defect $Q-W$:
\[
 \|Q-W\|\gtrsim U^{-m_h-1}.
\]
This does not rule out $Q\to W$, because the lower scale itself tends to zero; it rules
out only faster norm convergence.

Proposition~\ref{prop:o3r-polar-countermodel} shows that the square-root witness belongs
to the modulus layer, not the polar layer.  No lower bound for
$(I-P)U_SW$, no control of $U_R$ or $U_S$, and no conclusion about
$\Gamma_U\to I$, $K_{R,S}^{T_0,U}\to I$, strong terminal transport, Object X, Seal, or
RH follows.  The next genuine gate is an identity or estimate that couples the modulus
off-block $\mathscr N_U$ to the polar orientation of the concrete P11 metric pairs;
such a coupling is absent from abstract compression algebra.
\end{remark}

% P11 paper module: R20 individual noncommutativity versus relative polar incompatibility.

\subsection{Individual metric noncommutativity versus relative polar incompatibility}
\label{subsec:o3s-relative-polar}

Keep the odd-sector notation of Subsections~\ref{subsec:o3n-gauge-intertwining} and
\ref{subsec:o3r-sqrt-offblock}.  For a single source level $X\in\{R,S\}$ write
\[
 B_X:=G_{X,T_0}^-,\qquad C_X:=G_{X,U}^-,
\]
\[
 X_X=C_X^{1/2}B_X^{-1/2}=U_XH_X,
 \qquad
 H_X:=\bigl(B_X^{-1/2}C_XB_X^{-1/2}\bigr)^{1/2}.
\]
Proposition~\ref{prop:o3n-gauge-commutator} gives the qualitative equivalence
$U_X=I\iff[B_X,C_X]=0$.  The first point below is the scale-correct quantitative version
for one metric pair.

\begin{proposition}[Normalized square-root commutator controls individual polar activity]
\label{prop:o3s-individual-commutator}
For any boundedly invertible positive pair $B,C$, let
\[
 X=C^{1/2}B^{-1/2}=UH,
 \qquad
 H=(B^{-1/2}CB^{-1/2})^{1/2}.
\]
Then
\[
 X-X^*
 =B^{-1/2}[B^{1/2},C^{1/2}]B^{-1/2}
\tag{O3S.1}\label{eq:o3s-skew}
\]
and
\[
 \boxed{
 \|U-I\|
 \ge
 \frac{
 \|B^{-1/2}[B^{1/2},C^{1/2}]B^{-1/2}\|
 }{2\|H\|}.
 }
\tag{O3S.2}\label{eq:o3s-U-lower}
\]
The right-hand side is invariant under the common rescaling $(B,C)\mapsto(tB,tC)$,
$t>0$.
\end{proposition}

\begin{proof}
The identity~\eqref{eq:o3s-skew} is immediate from
\[
 B^{-1/2}[B^{1/2},C^{1/2}]B^{-1/2}
 =C^{1/2}B^{-1/2}-B^{-1/2}C^{1/2}.
\]
Since $X=UH$ and $X^*=HU^*$,
\[
 X-X^*=(U-I)H+H(I-U^*).
\]
Hence
\[
 \|X-X^*\|\le2\|H\|\,\|U-I\|,
\]
which proves~\eqref{eq:o3s-U-lower}.  Under common scalar rescaling both $X$ and $H$
are unchanged, and so is the normalized expression.
\end{proof}

The actual P11 gauge is not the individual quantity $U_X-I$ but the relative defect
\[
 \mathfrak P_U:=U_SW-WU_R.
\tag{O3S.3}\label{eq:o3s-relative-defect}
\]
Let $P=WW^*$ and retain
\[
 \Gamma_U=W^*U_SWU_R^*.
\]

\begin{proposition}[Exact relative polar decomposition]
\label{prop:o3s-relative-decomposition}
Put
\[
 \Lambda_U:=(I-P)U_SW.
\]
Then
\[
 \boxed{
 \mathfrak P_U
 =\Lambda_U+W(\Gamma_U-I)U_R,
 }
\tag{O3S.4}\label{eq:o3s-relative-decomp}
\]
and the two summands have orthogonal ranges.  Consequently, for every source vector
$f$,
\[
 \boxed{
 \|\mathfrak P_Uf\|^2
 =\|\Lambda_Uf\|^2
  +\|(\Gamma_U-I)U_Rf\|^2.
 }
\tag{O3S.5}\label{eq:o3s-relative-pythagoras}
\]
In particular a lower bound for the polar leakage $\Lambda_U$ would be sufficient for a
lower bound on the true relative gauge defect, but a lower bound on either individual
$\|U_R-I\|$ or $\|U_S-I\|$ is not.
\end{proposition}

\begin{proof}
Using $P=WW^*$ and the definition of $\Gamma_U$,
\[
\begin{aligned}
 \mathfrak P_U
 &=(I-P)U_SW+PU_SW-WU_R\\
 &=\Lambda_U+W(W^*U_SW-U_R)\\
 &=\Lambda_U+W(\Gamma_U-I)U_R.
\end{aligned}
\]
The first term lies in $(\Ran W)^\perp$ and the second in $\Ran W$, proving the
orthogonality and~\eqref{eq:o3s-relative-pythagoras}.
\end{proof}

The next model shows that individual metric noncommutativity, even if arbitrarily large
and even if it produces genuinely nontrivial polar factors, does not force relative
polar incompatibility.

\begin{proposition}[Canonical nested noncommutativity can cancel perfectly]
\label{prop:o3s-matched-countermodel}
There is a canonical inclusion model with exact base/future metric pullback in which
both source and target metric pairs are noncommuting, their raw commutator norms can be
made arbitrarily large, and $U_R,U_S$ are nontrivial, but
\[
 \boxed{
 (I-P)U_SW=0,
 \qquad
 U_SW-WU_R=0
 }
\tag{O3S.6}\label{eq:o3s-matched-zero}
\]
identically.
\end{proposition}

\begin{proof}
Let
\[
 B=\begin{pmatrix}1&0\\0&2\end{pmatrix},
 \qquad
 C=\begin{pmatrix}2&1\\1&2\end{pmatrix}.
\]
Then $B,C>0$ and
\[
 [B,C]=\begin{pmatrix}0&-1\\1&0\end{pmatrix}\ne0.
\]
Take source $\mathbb C^2$, target $\mathbb C^3$, and the canonical inclusion
$Wz=(z,0)$.  For $t>0$ set
\[
 G_{R,T_0}=tB,
 \qquad G_{R,U}=tC,
\]
\[
 G_{S,T_0}=tB\oplus1,
 \qquad G_{S,U}=tC\oplus1.
\]
The base and future metric pullback identities hold exactly.  Moreover the normalized
base transition is precisely the canonical inclusion $W$.  The polar products satisfy
\[
 X_S=X_R\oplus1,
\]
so uniqueness of polar decomposition gives
\[
 U_S=U_R\oplus1.
\]
Since $[B,C]\ne0$, Proposition~\ref{prop:o3n-gauge-commutator} gives $U_R\ne I$ and
$U_S\ne I$.  Nevertheless $U_SW=WU_R$, proving~\eqref{eq:o3s-matched-zero}.  Finally
\[
 \|[G_{R,T_0},G_{R,U}]\|
 =t^2\|[B,C]\|,
\]
and the same holds on the target nontrivial block, so the raw commutators are
arbitrarily large.
\end{proof}

The previous model has no modulus leakage.  The following direct-sum strengthening
shows that adding an R19-type modulus defect does not repair the implication.

\begin{proposition}[Modulus leakage plus large individual commutators still need not produce polar leakage]
\label{prop:o3s-combined-countermodel}
There is a canonical inclusion model satisfying exact base/future pullback such that
\[
 \|(I-P)Q\|>0,
\]
its square-root off-block can be made arbitrarily large, and both individual metric-pair
commutators can simultaneously be made arbitrarily large, while nevertheless
\[
 \boxed{
 (I-P)U_SW=0,
 \qquad
 U_SW-WU_R=0.
 }
\tag{O3S.7}\label{eq:o3s-combined-zero}
\]
\end{proposition}

\begin{proof}
For the first direct summand take the R19 canonical inclusion block
\[
 W_1:\mathbb C\hookrightarrow\mathbb C^2,
 \qquad W_1z=(z,0),
\]
with base metrics $1$ and $I_2$, and future metrics
\[
 2s^2,
 \qquad
 s^2A_0,
 \qquad
 A_0=\begin{pmatrix}2&1\\1&2\end{pmatrix}.
\]
Its polar factors are $1$ and $I_2$, while
\[
 Q_1=A_0^{1/2}W_1/\sqrt2
\]
has nonzero off-range component because $e_1$ is not an eigenvector of $A_0$; the
corresponding square-root off-block grows linearly in $s$.

For the second summand use the same noncommuting pair $B,C$ as in
Proposition~\ref{prop:o3s-matched-countermodel}, with identity inclusion between two
copies of $\mathbb C^2$, base metric $tB$ and future metric $tC$ on both copies.  Its
source and target polar factors are the same nontrivial unitary $U_0$, while the raw
commutator norm is $t^2\|[B,C]\|$.

Take the orthogonal direct sums of the two systems.  The total transition is a canonical
coordinate inclusion, the base/future pullback identities remain exact, and
\[
 U_R=1\oplus U_0,
 \qquad
 U_S=I_2\oplus U_0.
\]
Hence~\eqref{eq:o3s-combined-zero} holds.  At the same time the first summand supplies
nonzero modulus leakage and arbitrarily large square-root off-block as $s\to\infty$,
while the second supplies arbitrarily large individual commutators as $t\to\infty$.
\end{proof}

\begin{remark}[R20 firewall: the missing object is relative, not individual]
\label{rem:o3s-firewall}
Proposition~\ref{prop:o3s-individual-commutator} gives a legitimate quantitative route
from a scale-normalized square-root commutator to individual polar activity
$\|U_X-I\|$.  Propositions~\ref{prop:o3s-matched-countermodel} and
\ref{prop:o3s-combined-countermodel} show, however, that neither large raw
noncommutativity nor the combination of nontrivial modulus leakage with large
individual noncommutativity can be promoted by abstract algebra to a lower bound on
\[
 (I-P)U_SW
 \quad\text{or}\quad
 U_SW-WU_R.
\]
Thus the proposed target ``prove the two concrete metric pairs are quantitatively
noncommuting'' is insufficient as stated.  The remaining problem for this
polar-gauge obstruction route is a relative polar-compatibility problem: one must
detect a failure of the source and target polar rotations to intertwine under $W$.

Equivalently, the genuinely decisive quantity for this gauge component remains
\[
 \mathfrak P_U=U_SW-WU_R,
\]
or, by Proposition~\ref{prop:o3n-gauge-criterion}, the true gauge compression
$\Gamma_U-I$.  Any useful commutator criterion for this route must therefore be
relative across the nested metric pairs, not a pair of separate estimates on
$[G_{R,T_0},G_{R,U}]$ and $[G_{S,T_0},G_{S,U}]$.

No conclusion about $\Gamma_U\to I$, $K_{R,S}^{T_0,U}\to I$, strong terminal
transport, Object X, Seal, or RH follows here.
\end{remark}

% P11 paper module: R21 exact cross-polar defect and relative gauge reparametrization.

\subsection{Cross-polar defect and exact relative gauge reparametrization}
\label{subsec:o3t-crosspolar}

Keep the odd-sector notation at fixed $0<R<S<T_0<U$ from
Subsections~\ref{subsec:o3n-gauge-intertwining},
\ref{subsec:o3r-sqrt-offblock}, and~\ref{subsec:o3s-relative-polar}.  Thus
\[
 W:=W_{R,S,-}^{[T_0]},
 \qquad
 A_R:=A_{T_0,U}^{R,-},
 \qquad
 A_S:=A_{T_0,U}^{S,-},
\]
\[
 X_R=U_RA_R^{1/2},
 \qquad
 X_S=U_SA_S^{1/2},
\]
and the genuine relative polar defect is
\[
 \mathfrak P_U:=U_SW-WU_R.
\]
The purpose of this subsection is to express this defect through a single cross-product
of the two metric polar products.  No separate estimate of the individual commutators
$[G_{R,T_0},G_{R,U}]$ and $[G_{S,T_0},G_{S,U}]$ is used.

\begin{proposition}[Exact cross-polar identity]
\label{prop:o3t-crosspolar}
Define
\[
 \boxed{
 Z_U
 :=A_S^{-1/2}X_S^*WX_RA_R^{-1/2}.
 }
\]
Then
\[
 \boxed{
 Z_U=U_S^*WU_R.
 }
\tag{O3T.1}\label{eq:o3t-Z}
\]
In particular $Z_U$ is an isometry from the $R$-source space into the $S$-source
space.
\end{proposition}

\begin{proof}
Since $X_S^*=A_S^{1/2}U_S^*$ and $X_R=U_RA_R^{1/2}$,
\[
\begin{aligned}
 Z_U
 &=A_S^{-1/2}(A_S^{1/2}U_S^*)W(U_RA_R^{1/2})A_R^{-1/2}\\
 &=U_S^*WU_R.
\end{aligned}
\]
Both polar factors are unitary and $W$ is an isometry, hence so is $Z_U$.
\end{proof}

\begin{proposition}[Exact isometric gauge-defect reparametrization]
\label{prop:o3t-gauge-defect}
One has the exact operator identity
\[
 \boxed{
 Z_U-W=-U_S^*\mathfrak P_U.
 }
\tag{O3T.2}\label{eq:o3t-gauge-defect}
\]
Consequently, for every source vector $f$,
\[
 \boxed{
 \|(Z_U-W)f\|=\|\mathfrak P_Uf\|,
 }
\tag{O3T.3}\label{eq:o3t-vector-norm}
\]
and therefore
\[
 \boxed{
 Z_U-W\xrightarrow[s]{}0
 \quad\Longleftrightarrow\quad
 \mathfrak P_U\xrightarrow[s]{}0.
 }
\tag{O3T.4}\label{eq:o3t-strong-equivalence}
\]
In operator norm,
\[
 \boxed{
 \|Z_U-W\|=\|\mathfrak P_U\|,
 }
\tag{O3T.5}\label{eq:o3t-norm-equality}
\]
so norm convergence is equivalent as well.
\end{proposition}

\begin{proof}
Equation~\eqref{eq:o3t-Z} gives
\[
 Z_U-W
 =U_S^*WU_R-U_S^*U_SW
 =-U_S^*(U_SW-WU_R)
 =-U_S^*\mathfrak P_U.
\]
Unitary invariance of the vector and operator norms yields
\eqref{eq:o3t-vector-norm} and~\eqref{eq:o3t-norm-equality}.  The pointwise norm
identity gives the strong-convergence equivalence~\eqref{eq:o3t-strong-equivalence}.
\end{proof}

The preceding proposition is an exact reparametrization rather than an estimate: no
information about the relative polar defect is lost in passing from $\mathfrak P_U$ to
$Z_U-W$.  The same information can be written before the inverse-square-root
normalization as follows.

\begin{corollary}[Unnormalised relative cross-polar defect]
\label{cor:o3t-Crel}
Put
\[
 \boxed{
 \mathcal C_U^{\rm rel}
 :=X_S^*WX_R-A_S^{1/2}WA_R^{1/2}.
 }
\tag{O3T.6}\label{eq:o3t-Crel}
\]
Then
\[
 \boxed{
 \mathcal C_U^{\rm rel}
 =A_S^{1/2}(Z_U-W)A_R^{1/2}.
 }
\tag{O3T.7}\label{eq:o3t-Crel-factor}
\]
Hence
\[
 \boxed{
 \|\mathfrak P_U\|
 \ge
 \frac{\|\mathcal C_U^{\rm rel}\|}
 {\sqrt{\|A_S\|}\sqrt{\|A_R\|}},
 }
\tag{O3T.8}\label{eq:o3t-Crel-lower}
\]
while
\[
 \boxed{
 \|\mathfrak P_U\|
 \le
 \|A_S^{-1/2}\|\,
 \|\mathcal C_U^{\rm rel}\|\,
 \|A_R^{-1/2}\|.
 }
\tag{O3T.9}\label{eq:o3t-Crel-upper}
\]
For the concrete P11 family, Lemma~\ref{lem:o3l-relative-upper} gives
$\|A_R\|,\|A_S\|\ll e^U/U$, and therefore
\[
 \boxed{
 \|\mathfrak P_U\|
 \gtrsim
 \frac{U}{e^U}\,\|\mathcal C_U^{\rm rel}\|.
 }
\tag{O3T.10}\label{eq:o3t-P11-weighted-lower}
\]
The constants here may depend on the fixed levels $R,S,T_0$.
\end{corollary}

\begin{proof}
Using~\eqref{eq:o3t-Z},
\[
 X_S^*WX_R
 =A_S^{1/2}Z_UA_R^{1/2},
\]
which proves~\eqref{eq:o3t-Crel-factor}.  The first norm inequality follows from
\[
 \|\mathcal C_U^{\rm rel}\|
 \le
 \sqrt{\|A_S\|}\,\|Z_U-W\|\,\sqrt{\|A_R\|}
\]
and~\eqref{eq:o3t-norm-equality}.  Conversely,
\[
 Z_U-W
 =A_S^{-1/2}\mathcal C_U^{\rm rel}A_R^{-1/2},
\]
which gives~\eqref{eq:o3t-Crel-upper}.  Finally insert the relative-metric upper bounds
from Lemma~\ref{lem:o3l-relative-upper} into~\eqref{eq:o3t-Crel-lower}.
\end{proof}

There is also a source-space compression which is unitarily conjugate to the R15 gauge
compression.

\begin{corollary}[Cross-polar compression of the gauge]
\label{cor:o3t-Omega}
Define
\[
 \boxed{
 \Omega_U:=Z_U^*W
 =A_R^{-1/2}X_R^*W^*X_SA_S^{-1/2}W.
 }
\]
Then
\[
 \boxed{
 \Omega_U=U_R^*\Gamma_UU_R,
 }
\tag{O3T.11}\label{eq:o3t-Omega}
\]
where $\Gamma_U=W^*U_SWU_R^*$.  Consequently
\[
 \boxed{
 \|\Omega_U-I\|=\|\Gamma_U-I\|.
 }
\tag{O3T.12}\label{eq:o3t-Omega-norm}
\]
\end{corollary}

\begin{proof}
By~\eqref{eq:o3t-Z},
\[
 \Omega_U
 =U_R^*W^*U_SW
 =U_R^*(W^*U_SWU_R^*)U_R
 =U_R^*\Gamma_UU_R.
\]
The operator-norm identity follows by unitary invariance.
\end{proof}

\begin{remark}[No automatic strong-convergence transfer through $\Omega_U$]
\label{rem:o3t-Omega-strong-firewall}
Equation~\eqref{eq:o3t-Omega} is an exact unitary conjugacy at each fixed $U$, and
\eqref{eq:o3t-Omega-norm} is therefore exact.  Since the conjugating unitary $U_R$
depends on $U$, one must not infer from this identity alone an equivalence between
strong convergence of $\Omega_U$ and strong convergence of $\Gamma_U$.  The strong
statement that is exact without an additional hypothesis is instead
\eqref{eq:o3t-strong-equivalence} for $Z_U-W$ and $\mathfrak P_U$.
\end{remark}

\begin{remark}[R21-D scaling firewall]
\label{rem:o3t-scale-firewall}
The raw size of $\mathcal C_U^{\rm rel}$ is not itself a scale-invariant gauge measure.
Indeed, at fixed base metrics, simultaneously rescale the two future metrics by the
same scalar $t>0$.  Then
\[
 A_R,A_S\mapsto tA_R,tA_S,
 \qquad
 X_R,X_S\mapsto\sqrt t\,X_R,\sqrt t\,X_S,
\]
while the polar factors $U_R,U_S$, hence $Z_U-W$ and $\mathfrak P_U$, are unchanged.
In contrast,
\[
 \mathcal C_U^{\rm rel}\mapsto t\mathcal C_U^{\rm rel}.
\]
Thus a large unnormalised cross-polar norm cannot be promoted to a large gauge defect
without the metric normalization in~\eqref{eq:o3t-Crel-lower}--\eqref{eq:o3t-Crel-upper}.
\end{remark}

\begin{remark}[R21 countermodel firewall]
\label{rem:o3t-countermodels}
The cross-polar defect passes the canonical inclusion countermodels that blocked the
previous routes:
\begin{center}
\begin{tabular}{@{}p{0.23\textwidth}p{0.46\textwidth}p{0.21\textwidth}@{}}
\toprule
Model & Existing phenomenon & Cross-polar verdict \\
\midrule
R14, Proposition~\ref{prop:o3m-countermodel}
& $Q=W$ and $\Theta=0$, but the actual polar gauge is nontrivial
& $Z_U-W\ne0$ \\
R19, Proposition~\ref{prop:o3r-polar-countermodel}
& nonzero/arbitrarily large modulus and square-root leakage, but $U_S=I$, $U_R=I$
& $Z_U-W=0$ \\
R20 matched, Proposition~\ref{prop:o3s-matched-countermodel}
& nontrivial individual polar factors and arbitrarily large raw commutators, but $U_SW=WU_R$
& $Z_U-W=0$ \\
R20 combined, Proposition~\ref{prop:o3s-combined-countermodel}
& modulus leakage plus large individual commutators, still $U_SW=WU_R$
& $Z_U-W=0$ \\
\bottomrule
\end{tabular}
\end{center}
These conclusions are not separate model calculations: they follow immediately from
the exact identity~\eqref{eq:o3t-gauge-defect}.  Thus $Z_U-W$ rejects precisely the
matched polar configurations that should vanish and detects the R14 gauge-active
configuration.
\end{remark}

\begin{openproblem}[R21-F: concrete P11 cross-polar asymptotics]
\label{open:o3t-crosspolar}
Determine the asymptotic behavior of
\[
 Z_U-W
 =A_S^{-1/2}X_S^*WX_RA_R^{-1/2}-W
\]
for the concrete P11 metric family.  For the strong terminal problem the relevant task
is fixed-vector control of $(Z_U-W)f$; an operator-norm lower bound alone would not rule
out strong convergence.  Conversely, a norm upper bound tending to zero would imply
strong gauge compatibility.  Equivalently one may estimate
$\mathcal C_U^{\rm rel}$ together with sufficiently sharp control of the positive and
inverse square roots of $A_R,A_S$.

The present fixed-pair rank-one asymptotics do not provide this control:
Proposition~\ref{prop:o3n-rankone-insufficiency} shows that the inverse-square-root scale
is not determined by the existing leading Gram data, and Remark~\ref{rem:o3p-firewall}
identifies the remaining problem as operator-valued orientation of positive/inverse
square roots under $R\to S$.  Hence no conclusion about
\[
 \Gamma_U\to I,
 \qquad
 K_{R,S}^{T_0,U}\to I,
 \qquad
 W_{R,S,-}^{[U]}\text{ being strong Cauchy},
\]
or the existence of a global Object~X is drawn here.
\end{openproblem}

% P11 paper module: R22 positive strong-gauge angle defect and fixed-vector criterion.

\subsection{Positive strong-gauge angle defect and fixed-vector criterion}
\label{subsec:o3u-angle-defect}

Keep the odd-sector notation at fixed $0<R<S<T_0<U$ from
Subsections~\ref{subsec:o3n-gauge-intertwining},
\ref{subsec:o3s-relative-polar}, and~\ref{subsec:o3t-crosspolar}.  Thus
\[
 W:=W_{R,S,-}^{[T_0]},
 \qquad
 V_U:=U_SWU_R^*,
 \qquad
 \Gamma_U:=W^*V_U,
\]
\[
 \mathfrak P_U:=U_SW-WU_R,
 \qquad
 Z_U:=U_S^*WU_R,
 \qquad
 \Omega_U:=Z_U^*W=U_R^*\Gamma_UU_R.
\]
R15 identifies $\Gamma_U\to I$ strongly with $V_U\to W$ strongly.  The purpose of
this subsection is to encode that strong gate by a positive uniformly bounded
fixed-vector observable.  The moving-gauge operator built from $Z_U-W$ is retained,
but is not promoted to the R15 strong gate.

\begin{proposition}[Positive angle defect for the R15 strong gauge]
\label{prop:o3u-true-angle}
Define
\[
 \boxed{
 \mathscr G_U:=(V_U-W)^*(V_U-W).
 }
\tag{O3U.1}\label{eq:o3u-G-def}
\]
Then
\[
 \boxed{
 \mathscr G_U=2I-\Gamma_U-\Gamma_U^*.
 }
\tag{O3U.2}\label{eq:o3u-G-gamma}
\]
Moreover
\[
 \boxed{0\le\mathscr G_U\le4I,}
\tag{O3U.3}\label{eq:o3u-G-bound}
\]
and for every fixed source vector $f$,
\[
 \boxed{
 \langle\mathscr G_Uf,f\rangle
 =\|(V_U-W)f\|^2.
 }
\tag{O3U.4}\label{eq:o3u-G-rayleigh}
\]
\end{proposition}

\begin{proof}
Since $V_U$ and $W$ are isometries,
\[
\begin{aligned}
 (V_U-W)^*(V_U-W)
 &=V_U^*V_U-V_U^*W-W^*V_U+W^*W\\
 &=2I-\Gamma_U^*-\Gamma_U,
\end{aligned}
\]
which proves~\eqref{eq:o3u-G-gamma}.  Positivity is immediate from
\eqref{eq:o3u-G-def}.  Also $\|V_U-W\|\le2$, hence
$\|\mathscr G_U\|\le4$ and therefore~\eqref{eq:o3u-G-bound}.  Equation
\eqref{eq:o3u-G-rayleigh} follows from the definition.
\end{proof}

\begin{proposition}[Dense-core criterion for strong gauge compatibility]
\label{prop:o3u-dense-core}
Let $\mathcal C$ be any fixed dense subspace of the $R$-source Hilbert space.  Then
\[
 \boxed{
 \Gamma_U\xrightarrow[s]{}I
 \quad\Longleftrightarrow\quad
 \langle\mathscr G_Uf,f\rangle\longrightarrow0
 \quad\text{for every }f\in\mathcal C.
 }
\tag{O3U.5}\label{eq:o3u-dense-core}
\]
In particular, failure of the R15 strong gauge condition is certified by one fixed
vector and one subsequence whenever
\[
 \boxed{
 \limsup_{n\to\infty}
 \langle\mathscr G_{U_n}f,f\rangle>0
 }
\tag{O3U.6}\label{eq:o3u-fixed-witness}
\]
for some fixed $f$ and $U_n\to\infty$.
\end{proposition}

\begin{proof}
By Proposition~\ref{prop:o3n-gauge-criterion},
\[
 \Gamma_U\xrightarrow[s]{}I
 \quad\Longleftrightarrow\quad
 V_U\xrightarrow[s]{}W.
\]
The forward implication in~\eqref{eq:o3u-dense-core} follows from
\eqref{eq:o3u-G-rayleigh}.  Conversely suppose the quadratic forms tend to zero on
$\mathcal C$.  Then
\[
 \|(V_U-W)f\|\to0
 \qquad(f\in\mathcal C).
\]
For arbitrary $x$ and $f\in\mathcal C$,
\[
 \|(V_U-W)x\|
 \le2\|x-f\|+\|(V_U-W)f\|,
\]
because $\|V_U-W\|\le2$ uniformly.  Taking $U\to\infty$ and then using density of
$\mathcal C$ proves strong convergence on the whole source space.  Statement
\eqref{eq:o3u-fixed-witness} is the contrapositive fixed-vector obstruction.
\end{proof}

\begin{remark}[Smooth odd source core]
\label{rem:o3u-smooth-core}
The previously established TC0 graph-core reduction supplies
$C_c^\infty((-R,R))_{\rm odd}$ as a dense core of the odd P11 source space.  Hence the
criterion~\eqref{eq:o3u-dense-core} may be tested on fixed smooth odd source vectors;
no $U$-dependent maximizing vector is needed for a positive strong-convergence proof.
Conversely, an operator-norm lower bound carried only by $U$-dependent directions does
not by itself disprove strong convergence.
\end{remark}

The cross-polar defect of O3t produces a second positive operator.  It is norm-equivalent
to $\mathscr G_U$ but is expressed in a moving source gauge.

\begin{proposition}[Moving-gauge angle defect]
\label{prop:o3u-moving-angle}
Define
\[
 \boxed{
 D_U:=(Z_U-W)^*(Z_U-W).
 }
\tag{O3U.7}\label{eq:o3u-D-def}
\]
Then
\[
 \boxed{
 D_U=\mathfrak P_U^*\mathfrak P_U
 =2I-\Omega_U-\Omega_U^*
 =U_R^*\mathscr G_UU_R.
 }
\tag{O3U.8}\label{eq:o3u-D-identities}
\]
Consequently
\[
 \boxed{
 \|D_U\|
 =\|\mathscr G_U\|
 =\|\mathfrak P_U\|^2
 =\|V_U-W\|^2.
 }
\tag{O3U.9}\label{eq:o3u-D-norm}
\]
If
\[
 \Lambda_U:=(I-WW^*)U_SW,
\]
then the R20 orthogonal decomposition also gives the operator identity
\[
 \boxed{
 D_U
 =\Lambda_U^*\Lambda_U
 +U_R^*(\Gamma_U-I)^*(\Gamma_U-I)U_R.
 }
\tag{O3U.10}\label{eq:o3u-D-positive-split}
\]
\end{proposition}

\begin{proof}
O3t gives $Z_U-W=-U_S^*\mathfrak P_U$, hence
\[
 D_U=\mathfrak P_U^*\mathfrak P_U.
\]
Also
\[
 V_U-W=\mathfrak P_UU_R^*,
\]
so
\[
 \mathscr G_U
 =U_R\mathfrak P_U^*\mathfrak P_UU_R^*,
\]
which proves $D_U=U_R^*\mathscr G_UU_R$.  Substituting
$\Omega_U=U_R^*\Gamma_UU_R$ into~\eqref{eq:o3u-G-gamma} gives
$D_U=2I-\Omega_U-\Omega_U^*$.  Equation~\eqref{eq:o3u-D-norm} follows from unitary
invariance and
$\|T^*T\|=\|T\|^2$.

Finally Proposition~\ref{prop:o3s-relative-decomposition} gives
\[
 \mathfrak P_U
 =\Lambda_U+W(\Gamma_U-I)U_R
\]
with orthogonal ranges.  Taking the product with its adjoint on the source side removes
the cross terms and yields~\eqref{eq:o3u-D-positive-split}.
\end{proof}

\begin{proposition}[Moving-unitary conjugacy does not preserve fixed-vector strong asymptotics]
\label{prop:o3u-moving-countermodel}
There is an abstract gauge system with fixed isometry $W$ for which
\[
 D_n\xrightarrow[s]{}0
\]
while
\[
 \mathscr G_n\not\xrightarrow[s]{}0.
\]
Thus the pointwise strong behavior of $D_U$ cannot be promoted by abstract unitary
conjugacy to the R15 strong gauge condition.
\end{proposition}

\begin{proof}
Let the source and target be $\ell^2(\mathbb N)$ and take $W=I$.  Put
\[
 V:=I-2P_{e_1},
\]
the reflection in the first coordinate.  Then
\[
 \mathscr G_n:=(V-I)^*(V-I)=4P_{e_1}
\]
for every $n$, so $\mathscr G_n$ does not converge strongly to zero.

Choose unitaries $U_{R,n}$ with
\[
 U_{R,n}^*e_1=e_n
\]
and set
\[
 U_{S,n}:=VU_{R,n}.
\]
Then
\[
 V_n=U_{S,n}WU_{R,n}^*=V
\]
for every $n$, while
\[
 \mathfrak P_n
 =U_{S,n}W-WU_{R,n}
 =(V-I)U_{R,n}.
\]
Therefore
\[
 D_n
 =\mathfrak P_n^*\mathfrak P_n
 =U_{R,n}^*(4P_{e_1})U_{R,n}
 =4P_{e_n}
 \xrightarrow[s]{}0.
\]
This proves the claimed non-promotion.  The example is a topology firewall only; it is
not a counterexample to the concrete P11 metric family.
\end{proof}

\begin{remark}[R22-D strong-topology firewall]
\label{rem:o3u-moving-firewall}
At each fixed $U$, $D_U$ and $\mathscr G_U$ are unitarily conjugate and therefore have
the same norm and spectrum.  Since the conjugating unitary $U_R$ depends on $U$, this
pointwise equivalence does not identify their fixed-vector strong limits.  Thus
$D_U$ is a legitimate norm diagnostic and an exact positive form of the O3t cross-polar
defect, but the primary fixed-vector observable for the R15 strong gate is
$\mathscr G_U$.
\end{remark}

The true angle defect also admits a direct metric-product representation.

\begin{proposition}[Metric representation of the strong-gauge angle defect]
\label{prop:o3u-metric-representation}
Using
\[
 U_S=X_SA_S^{-1/2},
 \qquad
 U_R^*=A_R^{-1/2}X_R^*,
\]
one has
\[
 \boxed{
 V_U=X_SA_S^{-1/2}WA_R^{-1/2}X_R^*,
 }
\tag{O3U.11}\label{eq:o3u-V-metric}
\]
and
\[
 \boxed{
 \Gamma_U
 =W^*X_SA_S^{-1/2}WA_R^{-1/2}X_R^*.
 }
\tag{O3U.12}\label{eq:o3u-Gamma-metric}
\]
Consequently
\[
 \boxed{
 \begin{aligned}
 \mathscr G_U
 ={}&2I
 -W^*X_SA_S^{-1/2}WA_R^{-1/2}X_R^*\\
 &-X_RA_R^{-1/2}W^*A_S^{-1/2}X_S^*W.
 \end{aligned}
 }
\tag{O3U.13}\label{eq:o3u-G-metric}
\]
\end{proposition}

\begin{proof}
The first identity is the polar formulas inserted into
$V_U=U_SWU_R^*$.  Left compression by $W^*$ gives~\eqref{eq:o3u-Gamma-metric}; taking
the adjoint and substituting into~\eqref{eq:o3u-G-gamma} gives
\eqref{eq:o3u-G-metric}.
\end{proof}

\begin{remark}[Inverse-square-root information firewall]
\label{rem:o3u-inverse-root-firewall}
Equation~\eqref{eq:o3u-G-metric} does not bypass the R17--R21 information barrier.
Direct control of $\langle\mathscr G_Uf,f\rangle$ for a fixed vector $f$ requires the
relative orientation of the factors $A_R^{-1/2}$ and $A_S^{-1/2}$ inside the cross
product.  Equivalently, evaluating $V_Uf$ requires control of
\[
 A_R^{-1/2}X_R^*f
 \quad\text{and then}\quad
 X_SA_S^{-1/2}W A_R^{-1/2}X_R^*f.
\]
The existing rank-one fixed-pair Gram asymptotics do not determine this inverse-square-root
orientation; Proposition~\ref{prop:o3n-rankone-insufficiency} and the O3t open problem
remain in force.
\end{remark}

\begin{openproblem}[R22-F: fixed-vector P11 strong-gauge asymptotics]
\label{open:o3u-angle-defect}
On the fixed smooth odd source core, determine whether
\[
 \boxed{
 \langle\mathscr G_Uf,f\rangle
 =\|(V_U-W)f\|^2
 \longrightarrow0
 }
\tag{O3U.14}\label{eq:o3u-open-target}
\]
for every fixed $f$, or construct a fixed $f$ and a subsequence $U_n\to\infty$ for
which the limsup is positive.  By Proposition~\ref{prop:o3u-dense-core}, the first
alternative is exactly the R15 strong polar-gauge condition
$\Gamma_U\to I$.

This open problem concerns only the polar gauge component.  Even if it is resolved
positively, the full future transport
\[
 W_{R,S,-}^{[U]}=U_SQU_R^*
\]
also contains the modulus comparison $Q-W$.  No strong terminal-transport conclusion
and no global Object~X construction is asserted here.
\end{openproblem}

% P11 paper module: R23 rank-one polar drift under perfect modulus coherence.

\subsection{Rank-one polar drift under perfect modulus coherence}
\label{subsec:o3v-rankone-polar-drift}

R22 isolates the fixed-vector strong-gauge observable
\(
\mathscr G_U=(V_U-W)^*(V_U-W)
\).
The next model strengthens the R14 modulus-only firewall in a direction suggested by
the actual P11 asymptotics: even a single monotonically growing rank-one future channel
can generate a persistent polar rotation while the modulus comparison is exact.

\begin{proposition}[Canonical one-source rank-one family]
\label{prop:o3v-rankone-family}
Let \(\mathcal H\) be finite-dimensional, let \(B>0\), and let
\(e\in\mathcal H\) be a unit vector.  Put
\[
 \beta:=\langle Be,e\rangle,
 \qquad
 J:\mathbb C\to\mathcal H,
 \qquad Jz:=ze,
\]
and use the base metrics
\[
 B_R:=\beta,
 \qquad B_S:=B.
\]
The normalized base inclusion is
\[
 W=B^{1/2}J\,\beta^{-1/2},
 \qquad
 w:=We_\mathbb C=\frac{B^{1/2}e}{\sqrt\beta}.
\]
For \(t\ge1\), define
\[
 A_R(t):=t,
 \qquad
 A_S(t):=I+(t-1)ww^*,
\tag{O3V.1}\label{eq:o3v-relative}
\]
and future metrics
\[
 C_R(t):=t\beta,
 \qquad
 C_S(t):=B^{1/2}A_S(t)B^{1/2}.
\tag{O3V.2}\label{eq:o3v-future}
\]
Then the future pullback is exact,
\[
 \boxed{J^*C_S(t)J=C_R(t),}
\tag{O3V.3}\label{eq:o3v-pullback}
\]
and the modulus isometry and Jensen square-root defect satisfy
\[
 \boxed{Q_t=W,\qquad \Theta_t=0}
\tag{O3V.4}\label{eq:o3v-perfect-modulus}
\]
for every \(t\ge1\).
Moreover, if
\[
 u:=\frac{Be}{\sqrt\beta},
\]
then
\[
 \boxed{C_S(t)=B+(t-1)uu^*,}
\tag{O3V.5}\label{eq:o3v-rankone-growth}
\]
so the future target metric is a fixed positive floor plus monotone rank-one growth.
\end{proposition}

\begin{proof}
Since \(\|w\|=1\),
\[
 A_S(t)w=tw,
 \qquad
 A_S(t)^{1/2}w=\sqrt t\,w.
\]
Also
\[
 B^{1/2}w=\frac{Be}{\sqrt\beta}=u,
\]
which gives~\eqref{eq:o3v-rankone-growth}.  Furthermore
\[
 \langle u,e\rangle
 =\frac{\langle Be,e\rangle}{\sqrt\beta}
 =\sqrt\beta,
\]
so
\[
 J^*C_S(t)J
 =\beta+(t-1)|\langle u,e\rangle|^2
 =t\beta=C_R(t).
\]
The relative compression is
\[
 W^*A_S(t)W=t=A_R(t),
\]
and therefore
\[
 Q_t
 =A_S(t)^{1/2}WA_R(t)^{-1/2}
 =W.
\]
Finally
\[
 W^*A_S(t)^{1/2}W=\sqrt t=A_R(t)^{1/2},
\]
which is exactly \(\Theta_t=0\).
\end{proof}

\begin{theorem}[Rank-one polar-drift limit]
\label{thm:o3v-polar-drift}
Let \(W_t\) be the actual future normalized inclusion associated with
Proposition~\ref{prop:o3v-rankone-family}:
\[
 W_t
 :=C_S(t)^{1/2}J\,C_R(t)^{-1/2}.
\]
Then
\[
 \boxed{
 W_t1
 \longrightarrow
 \frac{Be}{\|Be\|}.
 }
\tag{O3V.6}\label{eq:o3v-future-limit}
\]
The fixed baseline direction is
\[
 W1=\frac{B^{1/2}e}{\sqrt{\langle Be,e\rangle}}.
\]
Hence
\[
 \boxed{
 W_t\longrightarrow W
 \quad\Longleftrightarrow\quad
 e\text{ is an eigenvector of }B.
 }
\tag{O3V.7}\label{eq:o3v-alignment-criterion}
\]
\end{theorem}

\begin{proof}
By~\eqref{eq:o3v-rankone-growth},
\[
 \frac1tC_S(t)
 =\frac1tB+\left(1-\frac1t\right)uu^*
 \longrightarrow uu^*
\]
in operator norm.  Norm continuity of the positive square root gives
\[
 \frac1{\sqrt t}C_S(t)^{1/2}
 =\left(\frac1tC_S(t)\right)^{1/2}
 \longrightarrow(uu^*)^{1/2}
 =\frac{uu^*}{\|u\|}.
\]
Therefore
\[
\begin{aligned}
 W_t1
 &=\frac1{\sqrt\beta}
   \left(\frac1tC_S(t)\right)^{1/2}e\\
 &\longrightarrow
 \frac1{\sqrt\beta}\frac{uu^*e}{\|u\|}
 =\frac{u}{\|u\|}
 =\frac{Be}{\|Be\|},
\end{aligned}
\]
proving~\eqref{eq:o3v-future-limit}.

If \(e\) is an eigenvector of \(B\), the two normalized directions coincide.
Conversely, equality of the two unit vectors implies
\[
 Be=cB^{1/2}e
\]
for some \(c>0\).  Multiplying by \(B^{-1/2}\) gives
\[
 B^{1/2}e=ce,
\]
so \(e\) is an eigenvector of \(B^{1/2}\), hence of \(B\).
\end{proof}

\begin{corollary}[Explicit limiting R22 angle defect]
\label{cor:o3v-angle-defect}
In the model above the source polar factor is the identity and
\(Q_t=W\), hence the actual future transport equals the polar-gauge transport
\(V_t\).  Consequently the R22 angle defect is scalar and
\[
 \boxed{
 \lim_{t\to\infty}\mathscr G_t
 =\delta_B(e)^2 I_{\mathbb C},
 }
\tag{O3V.8}\label{eq:o3v-G-limit}
\]
where
\[
 \boxed{
 \delta_B(e)^2
 =2-2\frac{\langle e,B^{3/2}e\rangle}
 {\sqrt{\langle e,B^2e\rangle\langle e,Be\rangle}}.
 }
\tag{O3V.9}\label{eq:o3v-half-full-angle}
\]
Moreover
\[
 \delta_B(e)=0
 \quad\Longleftrightarrow\quad
 e\text{ is an eigenvector of }B.
\]
\end{corollary}

\begin{proof}
The source product is
\[
 C_R(t)^{1/2}B_R^{-1/2}=\sqrt t>0,
\]
so \(U_R(t)=1\).  Equation~\eqref{eq:o3v-perfect-modulus} and the exact polar-gauge
factorization give
\[
 W_t=U_S(t)W=V_t.
\]
Thus
\[
 \mathscr G_t
 =(W_t-W)^*(W_t-W)
 =\|W_t1-W1\|^2I_{\mathbb C}.
\]
Insert Theorem~\ref{thm:o3v-polar-drift}.  Since
\[
 \|Be\|^2=\langle e,B^2e\rangle,
 \qquad
 \|B^{1/2}e\|^2=\langle e,Be\rangle,
\]
and
\[
 \langle Be,B^{1/2}e\rangle
 =\langle e,B^{3/2}e\rangle,
\]
one obtains~\eqref{eq:o3v-half-full-angle}.  The zero criterion is
Theorem~\ref{thm:o3v-polar-drift}.
\end{proof}

\begin{corollary}[Continuous strengthening of the R14 countermodel]
\label{cor:o3v-R14-strengthening}
Take
\[
 P=
 \begin{pmatrix}
 \sqrt3/2&1/2\\
 1/2&1
 \end{pmatrix},
 \qquad
 B=P^2,
 \qquad
 e=e_1.
\]
Then \(\langle Be,e\rangle=1\) and
\[
 W1=Pe_1
 =\binom{\sqrt3/2}{1/2}.
\]
For the continuous family
\[
 A_t=I+(t-1)(W1)(W1)^*,
 \qquad t\to\infty,
\]
one has \(Q_t=W\) and \(\Theta_t=0\) identically, but
\[
 \boxed{
 W_t1\longrightarrow
 \frac{
 \binom{1}{(2+\sqrt3)/4}
 }{
 \sqrt{1+((2+\sqrt3)/4)^2}
 }
 \ne W1.
 }
\tag{O3V.10}\label{eq:o3v-R14-limit}
\]
In particular
\[
 \boxed{
 \lim_{t\to\infty}\mathscr G_t
 =\left(
 2-\frac{2+5\sqrt3}{\sqrt{23+4\sqrt3}}
 \right)I_{\mathbb C}
 }
\tag{O3V.11}\label{eq:o3v-R14-angle}
\]
with
\[
 2-\frac{2+5\sqrt3}{\sqrt{23+4\sqrt3}}
 \approx0.0513796574>0.
\]
Thus no alternating sequence is needed: a monotone rank-one future blow-up with perfect
modulus coherence can converge to the wrong polar-gauge limit.
\end{corollary}

\begin{proof}
A direct multiplication gives
\[
 P^2e_1
 =\binom{1}{(2+\sqrt3)/4}.
\]
Theorem~\ref{thm:o3v-polar-drift} gives~\eqref{eq:o3v-R14-limit}; substituting the two
unit directions into Corollary~\ref{cor:o3v-angle-defect} gives
\eqref{eq:o3v-R14-angle}.
\end{proof}

\begin{remark}[R23 information firewall]
\label{rem:o3v-firewall}
Corollary~\ref{cor:o3v-R14-strengthening} is stronger than the earlier R14 alternating
countermodel in one specific sense: the future metric now has a fixed positive floor and
a single monotone rank-one channel, while the modulus comparison remains exact at every
parameter value.  Nevertheless the R22 angle defect approaches a positive constant.
Thus rank-one dominance is not by itself a polar-synchronization mechanism.

The obstruction is the mismatch between the baseline half-power direction
\[
 B^{1/2}e
\]
and the rank-one-dominant future full-power direction
\[
 Be.
\]
This model is not a counterexample to the concrete P11 metric family.  In P11 the
existing O3p/O3q results are fixed-pair or finite-block statements, and positive square
root does not commute with compression.  Therefore one must not insert a finite jet
Riesz vector into~\eqref{eq:o3v-half-full-angle} and call the result the full P11 gauge.
\end{remark}

\begin{openproblem}[R23-F: operator-level dominant-channel polar angle]
\label{open:o3v-dominant-angle}
Determine whether the concrete P11 future metric possesses a gauge-relevant
operator-level or reducing-block asymptotic of the schematic form
\[
 G_{X,U}=B_X+\tau_U r_Ur_U^*+E_U,
 \qquad \tau_U\to\infty,
\]
with remainder control strong enough to pass through the positive square root.  If so,
identify the resulting analogue of the half-power/full-power angle
\[
 \frac{B_Xr_U}{\|B_Xr_U\|}
 \quad\text{versus}\quad
 \frac{B_X^{1/2}r_U}{\|B_X^{1/2}r_U\|}
\]
and determine whether the corresponding relative \(R/S\) polar orientations
intertwine.

The present fixed-pair jet asymptotics do not supply this operator-level transfer, so no
claim about the concrete P11 limit of \(\mathscr G_U\), strong terminal transport, or a
global Object~X is made here.
\end{openproblem}

% P11 paper module: R24 uniform coercivity and resolvent bridge.

\subsection{Uniform coercivity and the inverse-root resolvent bridge}
\label{subsec:o3w-resolvent-bridge}

The R22 metric representation contains the nested inverse square roots
$A_R^{-1/2}$ and $A_S^{-1/2}$.  For the concrete P11 relative metrics their norms are,
however, uniformly controlled at fixed base terminal.  This separates the lower-spectral
question from the remaining orientation problem and gives an exact resolvent formula for
the latter.

\begin{proposition}[Uniform coercivity of the P11 relative metrics]
\label{prop:o3w-coercivity}
Fix $0<R<S<T_0$ and let $X\in\{R,S\}$.  For every $U>T_0$, put
\[
 A_X:=A_{T_0,U}^{X,-}.
\]
Then
\[
 \boxed{
 a_0 I\le A_X,
 \qquad
 a_0:=\frac1{1+\|H_{T_0}\|^2}>0.
 }
\tag{O3W.1}\label{eq:o3w-lower}
\]
Consequently
\[
 \boxed{
 \|A_X^{-1/2}\|
 \le a_0^{-1/2}
 =\sqrt{1+\|H_{T_0}\|^2}
 }
\tag{O3W.2}\label{eq:o3w-invroot-bound}
\]
uniformly in $U$.
Together with Lemma~\ref{lem:o3l-relative-upper},
\[
 a_0I\le A_X\le(1+\|H_U\|^2)I
 \ll\frac{e^U}{U}I.
\tag{O3W.3}\label{eq:o3w-two-sided}
\]
Thus the growth of the relative condition number comes from the upper spectral edge,
not from collapse of the lower edge.
\end{proposition}

\begin{proof}
For every nonzero source vector $z$, the Rayleigh quotient of $A_X$ is
\[
 \frac{q_U^X(E_{X,U}z)}{q_{T_0}^X(E_{X,T_0}z)}.
\]
Exact Gamma compatibility and positivity of the Schur term give
\[
 q_U^X(E_{X,U}z)\ge\mathfrak c_{\Gamma,X}[z].
\]
At the fixed terminal $T_0$, the graph-norm equivalence
\eqref{eq:graph-norm-equiv} gives
\[
 q_{T_0}^X(E_{X,T_0}z)
 \le(1+\|H_{T_0}\|^2)\mathfrak c_{\Gamma,X}[z].
\]
Their quotient is therefore at least $a_0$, proving~\eqref{eq:o3w-lower}.
Equation~\eqref{eq:o3w-invroot-bound} follows by functional calculus, and the upper
bound in~\eqref{eq:o3w-two-sided} is Lemma~\ref{lem:o3l-relative-upper}.
\end{proof}

\begin{remark}[Concrete inverse-root firewall]
\label{rem:o3w-concrete-firewall}
Proposition~\ref{prop:o3n-rankone-insufficiency} remains a valid abstract information
firewall for entrywise rank-one leading data, but the concrete P11 family has the extra
constraint~\eqref{eq:o3w-lower}.  Hence one must not interpret the remaining P11 gate as
possible blow-up of $\|A_X^{-1/2}\|$ caused by
$\lambda_{\min}(A_X)\downarrow0$.  The inverse-root norms are uniformly bounded.  What
remains unresolved is their orientation across the nested $R/S$ levels.
\end{remark}

Keep
\[
 A_R:=A_{T_0,U}^{R,-},
 \qquad
 A_S:=A_{T_0,U}^{S,-},
 \qquad
 W:=W_{R,S,-}^{[T_0]},
 \qquad
 P:=WW^*.
\]
The exact relative compression identity gives
\[
 A_R=W^*A_SW.
\]
Therefore
\[
 \boxed{
 \mathscr B_U
 :=(I-P)A_SW
 =A_SW-WA_R.
 }
\tag{O3W.4}\label{eq:o3w-B-intertwining}
\]
This is the same second-moment off-block used in O3l and O3r.

\begin{proposition}[Inverse-square-root resolvent bridge]
\label{prop:o3w-inverse-resolvent}
One has the exact operator-norm Bochner integral
\[
 \boxed{
 \begin{aligned}
 A_S^{-1/2}W-WA_R^{-1/2}
 ={}&-\frac1\pi\int_0^\infty t^{-1/2}
 (A_S+tI)^{-1}\mathscr B_U(A_R+tI)^{-1}\,dt.
 \end{aligned}
 }
\tag{O3W.5}\label{eq:o3w-inverse-resolvent}
\]
In particular
\[
 \boxed{
 \|A_S^{-1/2}W-WA_R^{-1/2}\|
 \le\frac1{2a_0^{3/2}}\|\mathscr B_U\|.
 }
\tag{O3W.6}\label{eq:o3w-inverse-bound}
\]
\end{proposition}

\begin{proof}
For a positive invertible operator $A$,
\[
 A^{-1/2}
 =\frac1\pi\int_0^\infty t^{-1/2}(A+tI)^{-1}\,dt.
\]
The resolvent quasi-commutator satisfies
\[
\begin{aligned}
 &(A_S+tI)^{-1}W-W(A_R+tI)^{-1}\\
 &\qquad
 =(A_S+tI)^{-1}
 \bigl(W(A_R+tI)-(A_S+tI)W\bigr)
 (A_R+tI)^{-1}\\
 &\qquad
 =-(A_S+tI)^{-1}\mathscr B_U(A_R+tI)^{-1},
\end{aligned}
\]
which proves~\eqref{eq:o3w-inverse-resolvent}.  By
Proposition~\ref{prop:o3w-coercivity},
\[
 \|(A_X+tI)^{-1}\|\le(a_0+t)^{-1}.
\]
Hence the integral converges in norm, and
\[
 \frac1\pi\int_0^\infty
 \frac{t^{-1/2}}{(a_0+t)^2}\,dt
 =\frac1{2a_0^{3/2}},
\]
which gives~\eqref{eq:o3w-inverse-bound}.
\end{proof}

\begin{proposition}[Square-root resolvent bridge]
\label{prop:o3w-square-resolvent}
The corresponding positive-square-root intertwining defect is
\[
 \boxed{
 \begin{aligned}
 A_S^{1/2}W-WA_R^{1/2}
 ={}&\frac1\pi\int_0^\infty t^{1/2}
 (A_S+tI)^{-1}\mathscr B_U(A_R+tI)^{-1}\,dt.
 \end{aligned}
 }
\tag{O3W.7}\label{eq:o3w-square-resolvent}
\]
Consequently
\[
 \boxed{
 \|A_S^{1/2}W-WA_R^{1/2}\|
 \le\frac1{2\sqrt{a_0}}\|\mathscr B_U\|.
 }
\tag{O3W.8}\label{eq:o3w-square-bound}
\]
\end{proposition}

\begin{proof}
Use
\[
 A^{1/2}
 =\frac1\pi\int_0^\infty
 t^{-1/2}A(A+tI)^{-1}\,dt
\]
and
\[
 A(A+tI)^{-1}=I-t(A+tI)^{-1}.
\]
Thus
\[
\begin{aligned}
 &A_S(A_S+tI)^{-1}W
 -WA_R(A_R+tI)^{-1}\\
 &\qquad
 =-t\bigl((A_S+tI)^{-1}W-W(A_R+tI)^{-1}\bigr)\\
 &\qquad
 =t(A_S+tI)^{-1}\mathscr B_U(A_R+tI)^{-1}.
\end{aligned}
\]
Integration gives~\eqref{eq:o3w-square-resolvent}.  Finally
\[
 \frac1\pi\int_0^\infty
 \frac{t^{1/2}}{(a_0+t)^2}\,dt
 =\frac1{2\sqrt{a_0}},
\]
proving~\eqref{eq:o3w-square-bound}.
\end{proof}

\begin{corollary}[Common resolvent source of the positive and inverse-root off-blocks]
\label{cor:o3w-common-source}
Applying $I-P$ to~\eqref{eq:o3w-square-resolvent} gives
\[
 \boxed{
 \mathscr C_U
 =(I-P)A_S^{1/2}W
 =\frac1\pi\int_0^\infty t^{1/2}
 (I-P)(A_S+tI)^{-1}\mathscr B_U(A_R+tI)^{-1}\,dt.
 }
\tag{O3W.9}\label{eq:o3w-C-resolvent}
\]
Likewise
\[
 \boxed{
 (I-P)A_S^{-1/2}W
 =-\frac1\pi\int_0^\infty t^{-1/2}
 (I-P)(A_S+tI)^{-1}\mathscr B_U(A_R+tI)^{-1}\,dt.
 }
\tag{O3W.10}\label{eq:o3w-inverse-offblock}
\]
Hence O3r's square-root witness and the inverse-root orientation entering O3t/O3u are
different spectral weightings of the same second-moment off-block.
\end{corollary}

\begin{remark}[R24 reduction]
\label{rem:o3w-reduction}
The raw norm lower bound for $\mathscr B_U$ from O3l does not decide the inverse-root
integral: the two resolvents may strongly suppress precisely the large directions.  But
R24 replaces the vague phrase ``inverse-square-root orientation'' by the explicit
resolvent-smoothed quantity
\[
 (A_S+tI)^{-1}\mathscr B_U(A_R+tI)^{-1}.
\]
This is a concrete analytic object built from operators already present in P11; no new
auxiliary gauge is introduced.
\end{remark}

\begin{openproblem}[R24-F: fixed-vector resolvent-weighted off-block asymptotics]
\label{open:o3w-resolvent-offblock}
For fixed smooth odd source vectors $f$, determine the asymptotics of
\[
 \int_0^\infty t^{-1/2}
 (A_S+tI)^{-1}\mathscr B_U(A_R+tI)^{-1}f\,dt
\]
and of the corresponding $t^{1/2}$-weighted integral as $U\to\infty$.
A vanishing result on a fixed dense core would give a genuine nested inverse-root
compatibility statement; a fixed-vector positive limsup would give an inverse-root
orientation obstruction.  The existing raw second-moment witness alone decides neither
alternative.
\end{openproblem}

